\documentclass[11pt,a4paper]{article}
\usepackage[T1]{fontenc}
\usepackage[utf8]{inputenc}
\usepackage{mathptmx}
\usepackage{amsmath,amssymb}
\usepackage[margin=2.6cm]{geometry}
\usepackage{booktabs}
\usepackage[activate=false]{microtype}
\usepackage{xcolor}
\definecolor{revP}{RGB}{20,60,170}
\newcommand{\revP}[1]{#1}
\definecolor{revQ}{RGB}{140,30,120}
\newcommand{\revQ}[1]{{\color{revQ}#1}}
\usepackage[colorlinks=true,linkcolor=blue!50!black,citecolor=blue!50!black,urlcolor=blue!50!black]{hyperref}
\usepackage{enumitem}
\newcommand{\code}[1]{\texttt{#1}}

\newcommand{\eps}{\varepsilon}
\newcommand{\R}{\mathbb{R}}
\newcommand{\C}{\mathbb{C}}
\newcommand{\Mcl}{M_{\mathrm{cl}}}
\newcommand{\Acl}{A_{\mathrm{cl}}}
\newcommand{\hp}{^{\mathsf H}}
\newcommand{\Rcl}{R_{\mathrm{cl}}}   % conjugate transpose
\newcommand{\tp}{^{\mathsf T}}   % plain transpose
\setlength{\parskip}{3pt}

\title{The algebra behind the oracle\\[3pt]
{\large Companion note: sparse $+$ low-rank structure in the $H_\infty$ value, its
gradients, and the spectral abscissa}}
\author{C.~Leclercq \and E.~Cl\'ement}
\date{September 7, 2026}

\begin{document}
\maketitle

\noindent
{\small [This is the consolidated edition of the note: four rounds of review questions (August 31 -- September 3) have been folded into the exposition, so nothing here is answered in the margin --- everything is introduced in order, and the decisions the review produced are collected in \S\ref{sec:todo}.  \revP{The
September 6 edition adds \S\ref{sec:proj}, a projection-based reduced model for the inner
loop, GL-validated and cavity-untested, with the matching to-do items.  \revQ{The September 7 edition credits the
optimise-and-enrich loop to Aliyev, Benner, Mengi and Voigt (2019), adds the two
mechanisms of Mengi (2022, 2023) for the watch list and the margin constraint, and
completes the citation of the disk certificate (Schwerdtner, Mengi and Voigt 2020);
these passages are set in this colour.}]}\par\smallskip

The optimisation note treated the oracle as a black box.  This note opens the box.  The
message can be stated up front: \emph{every} quantity the optimiser consumes --- the gain
$\sigma_1(\omega)$, its gradient with respect to \emph{all} design parameters at once, the
location and polish of peaks, the certified global maximum, the stability margin, its
certificate, and its gradient --- reduces to three primitive operations:
\begin{center}
\fbox{\parbox{0.9\textwidth}{
\textbf{(P1)} one sparse LU factorisation, per complex point $s$, of the inverse of the
\emph{plant resolvent} $R(s)=(sM-A)^{-1}$ --- i.e.\ of the sparse matrix $R(s)^{-1}=sM-A$
(the resolvent itself is never formed; ``applying $R$'' means triangular solves);\\
\textbf{(P2)} triangular solves with that factorisation (and its conjugate transpose);\\
\textbf{(P3)} dense arithmetic on objects of size $r$ or $2$ (the controller order and the
coupling rank).}}
\end{center}
Nothing in the pipeline ever factorises, stores, or differentiates a matrix of closed-loop
size.  This is what ``exploiting sparse $+$ low-rank'' means concretely, and it is the
property that ports to the cavity, where (P1) becomes a sparse multifrontal factorisation
of the linearised Navier--Stokes operator and (P2)--(P3) are unchanged.

%======================================================================
\section{The one structural fact}\label{sec:structure}
\label{sec:structure}

Plant ($n$ complex states, $M$ symmetric positive definite mass matrix, $A$ sparse):
\[
M\dot q = Aq + b\,u + d,\qquad y = c\hp q,
\]
with actuator $b\in\C^n$, sensor $c\in\C^n$ (Gaussians of the positions, or free fields
through a fixed smoothing map), and $d$ the full-field disturbance.  Controller (order
$r=8$; $A_c$ tridiagonal with rows $(d_c,e_c,f_c)$; input matrix $B_c$ \emph{fixed};
output row $n\tp$ and direct gain $K$ free where the case allows):
\[
\dot x_c = A_c x_c + B_c\, y,\qquad u = n\tp x_c + K\,y .
\]
Substituting $u$ and $y$, the closed-loop pencil on the stacked state
$(q,x_c)\in\C^{n+r}$ is
\[
\Mcl=\begin{bmatrix} M & 0\\ 0 & I_r\end{bmatrix},
\qquad
\Acl=
\underbrace{\begin{bmatrix} A & 0\\ 0 & A_c\end{bmatrix}}_{\text{sparse }\oplus\text{ small}}
\;+\;
\underbrace{\begin{bmatrix} b & 0\\ 0 & B_c\end{bmatrix}}_{U\in\C^{(n+r)\times 2}}
\underbrace{\begin{bmatrix} K c\hp & n\tp\\ c\hp & 0\end{bmatrix}}_{V\in\C^{2\times(n+r)}}.
\]
\begin{equation}
\boxed{\;\Acl \;=\; \mathrm{blkdiag}(A,\,A_c)\;+\;UV,\qquad \mathrm{rank}(UV)\le 2.\;}
\label{eq:structure}
\end{equation}
Every formula below is a consequence of \eqref{eq:structure}.

\paragraph{More actuators and sensors.}  The rank-2 structure is the $1{\times}1$ case of a
general fact.  With actuator matrix
$B_a\in\C^{n\times n_a}$, sensor matrix $C_s\in\C^{n\times n_s}$ (so $y=C_s\hp q\in\C^{n_s}$,
$u\in\C^{n_a}$) and a MIMO controller ($B_c\in\R^{r\times n_s}$, $N\in\C^{n_a\times r}$,
$K\in\C^{n_a\times n_s}$), the coupling term of \eqref{eq:structure} becomes
$U\in\C^{(n+r)\times(n_a+n_s)}$, $V\in\C^{(n_a+n_s)\times(n+r)}$:
$\mathrm{rank}(UV)\le n_a+n_s$ --- still tiny.  Everything upgrades by replacing scalars
with small matrices: the controller collapses to the $n_a\times n_s$ transfer
\emph{matrix} $\mathsf K(s)=K+N(sI-A_c)^{-1}B_c$; Sherman--Morrison becomes Woodbury with
an $(n_a{+}n_s)$-sized capacity matrix; the plant transfer is
$\mathsf G(s)=C_s\hp R(s)B_a\in\C^{n_s\times n_a}$ ($n_a$ solves on one factorisation);
and the winding scalar becomes the \emph{return-difference determinant},
$f(s)=\det\bigl(I-\mathsf K(s)\mathsf G(s)\bigr)$ ($n_s\times n_s$, dense-tiny), with the
factorisation \eqref{eq:factorised} and the count formula \eqref{eq:count} holding
verbatim.  Costs: the LU is unchanged; triangular solves multiply by $n_a$ (forward) and
$n_s$ (adjoint); gradients still come from one forward/adjoint \emph{pair of blocks} per
frequency.  The capability is scheduled (\S\ref{sec:todo}, item~7).

%======================================================================
\section{Every closed-loop solve is a plant solve}
\label{sec:solve}

Let $s\in\C$ (on the frequency axis, $s=i\omega$) and consider the closed-loop system
$(s\Mcl-\Acl)\bigl(\begin{smallmatrix}q\\x_c\end{smallmatrix}\bigr)
=\bigl(\begin{smallmatrix}r_q\\r_c\end{smallmatrix}\bigr)$.
Eliminate the controller states from the second block row:
$x_c=(sI-A_c)^{-1}\bigl(r_c+B_c\,c\hp q\bigr)$ --- an $r\times r$ dense solve, primitive
(P3).  Substituting into the first row collapses the whole controller to a single scalar,
its transfer function at $s$:
\begin{equation}
\boxed{\;k(s)\;=\;K + n\tp (sI-A_c)^{-1}B_c,\;}
\label{eq:ks}
\end{equation}
and the plant-level system becomes
$\bigl(R(s)^{-1} - k(s)\,b\,c\hp\bigr)\,q = \tilde r$, with $R(s)=(sM-A)^{-1}$ the plant
resolvent.  A rank-one perturbation of a factorised matrix is solved by the
Sherman--Morrison identity:
\begin{equation}
\boxed{\;
\bigl(R^{-1} - k\,bc\hp\bigr)^{-1}
= R + \frac{k\; Rb\; c\hp R}{\,1-k\,G(s)\,},
\qquad G(s) = c\hp R(s)\, b .
\;}
\label{eq:sherman}
\end{equation}
So one closed-loop solve costs: one sparse LU of $R(s)^{-1}$ (P1, done once per $s$ and
reused for everything), two or three triangular solves (P2: applying $R$ to $\tilde r$
and to $b$; $G=c\hp Rb$ comes free from the second), and scalars (P3).  Adjoint systems
$(\,\cdot\,)\hp x = y$ use the \emph{same} factorisation via transposed-conjugate
triangular solves --- no second LU.

Note the scalar in the denominator of \eqref{eq:sherman},
\[
f(s) \;=\; 1 - k(s)\,G(s),
\]
the \emph{return difference} of the SISO loop.  Keep it in mind: it reappears as the star
of the stability certificate in \S\ref{sec:winding}.  (Numerically, $|f|$ small means the
solve is ill-conditioned --- which happens exactly when $s$ approaches a closed-loop
eigenvalue, as \S\ref{sec:winding} makes precise.)

%======================================================================
\section{The $H_\infty$ value: $\sigma_1(\omega)$ as an eigenvalue}
\label{sec:sigma}

The performance channel is full-field: disturbance $d$ into the plant equation, output the
plant state $q$ measured in the energy norm $\|q\|_M^2=q\hp Mq$.  With
$\Pi=[\,I_n\;\;0\,]$ the extractor of plant states, the (weighted) transfer operator at
frequency $\omega$ is
\[
T(\omega)\;=\;M^{1/2}\,\Pi\,(i\omega \Mcl-\Acl)^{-1}\,\Pi\tp M^{1/2}
\;\in\;\C^{n\times n},
\qquad
\sigma_1(\omega)=\|T(\omega)\|_2 .
\]
We never form $T$; we need only its largest singular value, i.e.\ the largest eigenvalue
of the Hermitian positive semidefinite operator $T\hp T$.  Unwinding the square roots,
$\sigma_1^2$ is the largest eigenvalue of
\begin{equation}
\boxed{\;
W(\omega)\;=\;M\,\Rcl(\omega)\hp\,M\,\Rcl(\omega),
\qquad
\Rcl(\omega)=\Pi(i\omega\Mcl-\Acl)^{-1}\Pi\tp ,
\;}
\label{eq:W}
\end{equation}
understood as an operator on $\C^n$ (self-adjoint in the $M^{-1}$ inner product; the code
applies it verbatim: solve, multiply by $M$, adjoint-solve, multiply by $M$).  One
application of $W$ = two solves of \S\ref{sec:solve} on one factorisation.  The largest
eigenvalue is obtained by power/Lanczos iteration on $W$, \emph{warm-started} with the
converged eigenvector from the neighbouring frequency --- across an exchange set the
vectors vary slowly in $\omega$, and the warm start cuts the iteration count to a
handful.  The converged eigenvector $v$ (the worst disturbance shape) and its image
give the singular pair $(u,\sigma_1,v)$ of $T$ that the gradients need next.

Cost of one oracle value: one LU $+$ $O(10)$ triangular solves.  On the 1-D GL mesh the LU
is $O(n)$; at cavity scale it is the sparse factorisation of a shifted Navier--Stokes
operator --- expensive, which is why the outer method was shaped so that \emph{many
independent frequencies} are evaluated concurrently.

%======================================================================
\section{Gradients of $\sigma_1$: all parameters from one solve pair}
\label{sec:grads}

Fix $\omega$ and suppose $\sigma_1$ is a \emph{simple} singular value of $T$ (the solver
monitors the gap $\sigma_1-\sigma_2$; more on ties in \S\ref{sec:subgrad}).  With unit
singular vectors $u,v$ ($Tv=\sigma_1 u$), first-order perturbation theory gives the exact
directional derivative
\[
\delta\sigma_1 \;=\; \mathrm{Re}\,\bigl[\,u\hp\,(\delta T)\,v\,\bigr].
\]
Since only $\Acl$ depends on the design, $\delta\bigl((i\omega\Mcl-\Acl)^{-1}\bigr)
=(i\omega\Mcl-\Acl)^{-1}\,\delta\Acl\,(i\omega\Mcl-\Acl)^{-1}$, and defining the two
\emph{solved vectors} (one forward, one adjoint solve --- both on the factorisation
already computed for the value):
\[
\tilde v=(i\omega\Mcl-\Acl)^{-1}\Pi\tp M^{1/2}v \in\C^{n+r},
\qquad
\tilde u=(i\omega\Mcl-\Acl)^{-\mathsf H}\Pi\tp M^{1/2}u \in\C^{n+r},
\]
we obtain the master formula
\begin{equation}
\boxed{\;\delta\sigma_1 \;=\; \mathrm{Re}\,\bigl[\,\tilde u\hp\,(\delta \Acl)\,\tilde v\,\bigr].\;}
\label{eq:master}
\end{equation}
Everything now follows from the \emph{structure} \eqref{eq:structure} of $\Acl$: each
design parameter perturbs $\Acl$ in a known sparse or low-rank pattern, so
\eqref{eq:master} evaluates to inner products of pieces of $\tilde u,\tilde v$.  Write
$\tilde v=(\tilde v_p,\tilde v_c)$ and $\tilde u=(\tilde u_p,\tilde u_c)$ for the plant
and controller components.  Then (for a complex parameter $p$ we record the complex
sensitivity $\zeta_p$ defined by $\delta\sigma_1=\mathrm{Re}[\bar\zeta_p\,\delta p]$; the
real gradient with respect to $(\mathrm{Re}\,p,\mathrm{Im}\,p)$ is
$(\mathrm{Re}\,\zeta_p,\mathrm{Im}\,\zeta_p)$):

\begin{center}
\small
\begin{tabular}{lll}
\toprule
parameter & enters $\Acl$ as & sensitivity from \eqref{eq:master} \\
\midrule
gain $K$ & $\delta K\,[\,b;0\,][\,c\hp\;0\,]$ &
 $\bar\zeta_K=(\tilde u_p\hp b)\,(c\hp\tilde v_p)$ \\[2pt]
output row $n$ & $[\,b;0\,][\,0\;\;\delta n\tp\,]$ &
 $\bar\zeta_{n_j}=(\tilde u_p\hp b)\,(\tilde v_c)_j$ \\[2pt]
$A_c$ entries $(d_c,e_c,f_c)$ & $\delta(A_c)_{jk}$ in the controller block &
 $\bar\zeta_{jk}=(\tilde u_c)^*_j\,(\tilde v_c)_k$ \quad (tridiagonal $jk$ only) \\[2pt]
actuator $b$ & $\delta b\,[\,Kc\hp\tilde v_p + n\tp\tilde v_c\,]$ direction &
 $\nabla_b$-vector $\;\propto\; \overline{\tilde u_p}\,\cdot
 \bigl(Kc\hp\tilde v_p + n\tp\tilde v_c\bigr)$ \\[2pt]
sensor $c$ & $[\,\tilde u_p\hp bK + \tilde u_c\hp B_c\,]\,\delta c\hp$ direction &
 $\nabla_c$-vector $\;\propto\; \tilde v_p\,\cdot
 \overline{\bigl(K\,\tilde u_p\hp b + \tilde u_c\hp B_c\bigr)}$ \\[2pt]
position $x_a$ & chain rule through $b=g(x_a)$ &
 $\partial\sigma/\partial x_a = \mathrm{Re}\bigl[\,\zeta_b\hp\, g'(x_a)\,\bigr]$
 \ (similarly $x_s$) \\
\bottomrule
\end{tabular}
\end{center}

\noindent
Free fields are optimised through a fixed linear smoothing (Sobolev) map
$b=\Sigma\,\beta$; the chain rule multiplies $\nabla_b$ by the constant $\Sigma\hp$.
The point to internalise is the \textbf{cost}: the entire gradient --- all
$\mathrm{NZ}$ real parameters, whether $\mathrm{NZ}=2$ (a position case) or
$\mathrm{NZ}=1659$ (a two-field case) --- costs \emph{two} triangular solve sets beyond
the value, independent of $\mathrm{NZ}$.  This is the adjoint-state economy, and it is why
the field cases are affordable at all.

%======================================================================
\section{Frequency derivatives, peak location, and polishing}

The same master formula with $\theta$ frozen and $\omega$ varied
($\delta(i\omega\Mcl-\Acl) = i\,\delta\omega\,\Mcl$) gives
\[
\frac{\partial\sigma_1}{\partial\omega}
= \mathrm{Re}\bigl[\,\tilde u\hp(-i\Mcl)\tilde v\,\bigr]
= \mathrm{Im}\bigl[\,\tilde u\hp \Mcl\,\tilde v\,\bigr]:
\]
one inner product per frequency, no extra solve.  A peak is a zero of
$\partial_\omega\sigma_1$; Newton's iteration on this scalar (second derivative by
differencing the first, or one extra solve) polishes each located peak to $10^{-10}$ in a
few oracle calls, warm-started along the iteration.  Peak \emph{location} starts from the
exchange set and from a coarse scan; \emph{missed} peaks are the certificate's job:

%======================================================================
\section{Certifying the global maximum}\label{sec:cert}
\label{sec:cert}

The optimiser must never trust a locally found maximum (v6, \S5).  The certified question
is: \emph{given a level $\gamma$, does $\sigma_1(\omega)<\gamma$ hold for all $\omega$ in
the band?}  The classical equivalence that turns this analytic question into an algebraic
one: $\gamma$ is a singular value of $T(\omega)$ if and only if the structured pencil
\begin{equation}
\boxed{\;
\mathcal P_\gamma(\lambda)\;=\;
\begin{bmatrix} \; -A & \;\gamma^{-1}\,\hat B\hat B\hp \\[2pt]
 -\gamma^{-1}\,\hat C\hp\hat C & \;-A\hp \end{bmatrix}
-\lambda
\begin{bmatrix} -M & 0\\ 0 & M\hp\end{bmatrix}
\;}
\label{eq:pencil}
\end{equation}
(written here for the closed-loop realisation, $\hat B,\hat C$ the weighted input/output
maps of \S\ref{sec:sigma}) is singular at $\lambda=i\omega$.  Consequently: if the pencil
has \emph{no eigenvalue} on a segment of the imaginary axis, then $\sigma_1\ne\gamma$
throughout that segment, and a single sample below $\gamma$ certifies the whole segment
below $\gamma$.  The oracle therefore covers the band by \textbf{eigenvalue-free disks}:
at a test point $i\omega_0$, a structure-preserving shift-invert Arnoldi iteration finds
the pencil eigenvalue nearest to $i\omega_0$; its distance is a certified radius within
which the level is not crossed.  Each shift-invert solve is a bordered sparse system that
reduces, once again, to (P1)--(P3) --- plant factorisations plus rank-$O(1)$ algebra.
Adaptive placement of test points sweeps the band; wherever a disk cannot be certified,
the iteration has \emph{found} an axis-near eigenvalue, i.e.\ located a frequency where
$\sigma_1$ reaches $\gamma$ --- exactly the ``missed peak'' the exchange needs.  Bisection
on $\gamma$ (warm-started from the incumbent level) certifies the maximum itself to
relative tolerance $10^{-6}$ at a cost of $O(10^2)$ shift-invert solves per design.

%======================================================================
\section{Stability: value, certificate, and the return difference}
\label{sec:winding}

\subsection{The value}
The margin constraint is $\alpha(\theta)\le-\eps$ with
$\alpha=\max_j\mathrm{Re}\,\lambda_j(\Acl,\Mcl)$.  During optimisation the pipeline does
not recompute the whole spectrum; it \emph{tracks} the few dangerous eigenvalues by
shift-invert (inverse) iteration, shifting at each accepted step to the eigenvalue's
previous location --- each iteration a bordered solve on the plant factorisation.  At
benchmark scale a dense generalized eigensolve provides ground truth; at cavity scale it
will not exist, which is why a certificate is required:

\subsection{The certificate: count eigenvalues without computing them}
Take determinants in the block form of $s\Mcl-\Acl$ and eliminate the controller block
(Schur complement), then apply the matrix determinant lemma to the remaining rank-one
term --- the same two moves as \S\ref{sec:solve}, now on determinants:
\begin{align}
\det(s\Mcl-\Acl)
&= \det(sI-A_c)\,\det\!\bigl(R(s)^{-1}-k(s)\,bc\hp\bigr) \notag\\
&= \det(sI-A_c)\;\det R(s)^{-1}\;\bigl(1-k(s)\,G(s)\bigr) .
\label{eq:factorised}
\end{align}

The passage from $R\,bc\hp$ to the scalar $G$ is the \emph{matrix determinant lemma}, in two
lines.  Factor
out: $\det(R^{-1}-k\,bc\hp)=\det R^{-1}\cdot\det(I-k\,(Rb)\,c\hp)$.  Now
$(Rb)\,c\hp$ is a rank-one matrix; its eigenvalues are $c\hp(Rb)=G$ (eigenvector $Rb$)
and $0$ with multiplicity $n-1$, so
$\det(I-k\,(Rb)c\hp)=\prod_i(1-k\lambda_i)=1-k\,G$.  The row and column vectors trade
places inside the determinant exactly because a rank-one matrix's only nonzero
eigenvalue is the scalar ``row times column''.
\begin{equation}
\boxed{\;\det(s\Mcl-\Acl)\;=\;\det R(s)^{-1}\,\det(sI-A_c)\;f(s),
\qquad f(s)=1-k(s)G(s).\;}
\end{equation}
The closed-loop spectrum is the union of the plant spectrum, the (open) controller
spectrum, and the zeros of the scalar $f$ --- with the crucial feature that the poles of
$f$ are exactly the plant and controller eigenvalues, so the products cancel where they
must.

Why the poles of $f$ must sit exactly at the open-loop eigenvalues: count orders at one
point.  Let $\lambda_p$ be a simple \emph{plant} eigenvalue
that is \emph{not} a closed-loop eigenvalue.  The left side $\det(s\Mcl-\Acl)$ is a
polynomial, finite and nonzero at $\lambda_p$.  On the right, $\det R(s)^{-1}=\det(sM-A)$
vanishes to first order at $\lambda_p$ --- so $f$ must blow up to first order there to
keep the product finite, and indeed it does: $G(s)=c\hp R(s)b$ has a simple pole at every
plant eigenvalue (the resolvent blows up), hence so does $f=1-kG$.  The same argument at
controller eigenvalues, where $k(s)$ carries the pole.  Nothing is being arranged; the
factorisation \eqref{eq:factorised} is an identity between a polynomial and a product of
two vanishing factors and one blowing factor, and the orders must balance --- which is
precisely why the pole count $P_{\mathrm{plant}}+P_{\mathrm{ctrl}}$ appears in the count
formula \eqref{eq:count}: the argument principle counts $f$'s zeros \emph{minus} its
poles, and the poles are exactly the open-loop eigenvalues the two determinant factors
account for.  Now apply the argument principle to $f$ on the boundary of the half-plane
$\{\mathrm{Re}\,s>\ell\}$ (the vertical line $\mathrm{Re}\,s=\ell$ closed by a large arc):
zeros minus poles of $f$ inside equals the winding number of $f$ along the contour, hence
\begin{equation}
\boxed{\;
\#\{\lambda_{\mathrm{cl}}:\mathrm{Re}\,\lambda_{\mathrm{cl}}>\ell\}
\;=\; W\bigl[f;\ \mathrm{Re}\,s=\ell\bigr] \;+\; P_{\mathrm{plant}}(\ell)
\;+\; P_{\mathrm{ctrl}}(\ell),
\;}
\label{eq:count}
\end{equation}



\paragraph{The winding number, defined.}
\label{sec:winding-num}
$W[f;\mathrm{Re}\,s=\ell]$ is the winding number of the curve $f(s)$ around the origin as
$s$ traverses the boundary of the region $\{\mathrm{Re}\,s>\ell\}$ with the region kept on
the left: \emph{down} the line, from $\ell+iY$ to $\ell-iY$, closed by the large arc
through $\ell+Y$.  Equivalently
$W=\tfrac1{2\pi}[\arg f]_{\mathrm{contour}}=\tfrac1{2\pi i}\oint f'/f\,ds$, and with this
orientation the argument principle reads $Z-P=W$ over the region, which is how
\eqref{eq:count} is meant.

\paragraph{How it is computed: exact phase tracking, not quadrature.}
Numerical integration of $f'/f$ would carry a quadrature error, and an error of $0.51$ in
a count is a wrong theorem.  Instead the code tracks the phase: sample $f$ along the
segment, and between consecutive samples add the branch-consistent increment
$\Delta\arg\in(-\pi,\pi]$; the sum telescopes to the exact total phase change
\emph{provided no true increment ever exceeded $\pi$ unnoticed}.  Three mechanisms
enforce that proviso, each added after a measured miscount during validation:
(i)~\emph{mandatory resolution near the poles of $f$} --- near an open-loop eigenvalue
close to the line the phase rotates fastest, so the sampler is forced to a density scaled
by the distance to the nearest pole; (ii)~\emph{midpoint validation} --- every accepted
increment is re-checked at the interval's midpoint: the two half-increments must sum to
the whole, a cheap proof that no $\pm2\pi$ was skipped; (iii)~a \emph{mandatory knot
grid}, so that local adaptivity cannot glide over a feature between coarse seeds.  With
all three, the certificate reproduced the dense count on $104/104$ validation designs.

\paragraph{Parallel phase tracking.}  The tracked quantity is a sum of increments, and sums split: the
segment is cut at the mandatory knots, each inter-knot interval goes to a worker that
runs the serial predictor--validator march and returns its phase subtotal, and the
subtotals add exactly --- no phase continuity across cuts is needed, because no absolute
branch of $\arg f$ is ever chosen (the only absolute arguments are at $\pm Y$, where $f$
is confined near $1$ and the principal branch is unambiguous).  Two passes: knot values
first (one parallel round); then intervals whose endpoint Hermite estimate already
validates the chord increment are closed at zero extra cost, and the rest go to a work
queue with midpoint splitting.  The wall time is the longest interval --- a few tens of
serial samples, set by the hottest pole window --- not the total count.  The knot grid is $121$ points over the core band, seven around every hot pole, plus the
tail knots to $\pm Y$; shared-knot values are computed once in pass~1, so no duplicated
factorisations; the three hardening mechanisms act \emph{locally}, so they parallelise
unchanged; and roundoff cannot flip a count, since an increment would have to move by
$\pi$.  Honest status: \code{winding\_cert.py} as frozen is serial; the decomposition is
\S\ref{sec:todo}, item~4, to be measured at GL where the serial reference exists.

\paragraph{The tail bound, derived slowly.}
The segment must end somewhere; the question is where the phase provably stops moving.
Start from the Neumann series: for any square matrix $B$ and $|s|>\|B\|$,
\[
(sI-B)^{-1}=\frac1s\Bigl(I-\frac Bs\Bigr)^{-1}=\frac1s\sum_{k\ge0}\Bigl(\frac Bs\Bigr)^{k},
\qquad\text{so}\qquad
\|(sI-B)^{-1}\|\le\frac1{|s|}\cdot\frac1{1-\|B\|/|s|}=\frac1{|s|-\|B\|}.
\]
Apply it to the plant resolvent by factoring out the mass matrix,
$R(s)=(sM-A)^{-1}=(sI-M^{-1}A)^{-1}M^{-1}$: with
\[
\nu:=\|M^{-1}A\|
\quad\text{(any upper bound will do; a few power iterations on $M^{-1}A$, once per mesh)},
\]
and $\|M^{-1}\|=1/\lambda_{\min}(M)$
\quad(one Lanczos run on the symmetric positive definite $M$, once per mesh; for a lumped
mass matrix, the smallest diagonal entry),
\[
\|R(s)\|\;\le\;\frac{1}{\lambda_{\min}(M)\,(|s|-\nu)}
\qquad\Longrightarrow\qquad
|G(s)|=|c\hp R(s)b|\;\le\;\frac{\|c\|\,\|b\|}{\lambda_{\min}(M)\,(|s|-\nu)} .
\]

\paragraph{Computing the once-per-mesh constants.}
Start from $\lambda_{\min}(M)=1/\lambda_{\max}(M^{-1})$: ``Lanczos
on $M$'' in the form that actually converges quickly is shift-invert Lanczos with shift
$0$, i.e.\ the Lanczos (power) iteration on the operator $x\mapsto M^{-1}x$.  Concretely:
one sparse Cholesky of $M$ --- $M$ is SPD, tridiagonal-banded on the 1-D mesh, a
nested-dissection factor in 2-D --- then each Lanczos step is two triangular solves, and
the largest eigenvalue of $M^{-1}$ is extremal, so it converges in a few tens of steps.
Sparsity is exploited throughout: $M$ is touched only through its factor.  This is the one
place the pipeline factorises the mass matrix, and it is admissible because it happens
once per mesh, outside every loop, and the factor is discarded afterwards; the rule
``never factorise $M$'' (v2 report, \S1.3) was about the optimisation loop.  Now the
honest accounting, in three thrift levels.  \emph{(a) What the shipped code does}
(\code{winding\_cert.py}, \code{plant\_tail\_constants}): \code{eigsh(M, k=1,
which='SM')}, i.e.\ Lanczos on $M$ itself in regular mode, asking ARPACK for the
smallest-magnitude Ritz value --- no factorisation at all, and fine at $n=399$, but it is
the slow way (the small end of a mass-matrix spectrum is the clustered end, and
regular-mode Lanczos converges to a clustered extremal eigenvalue in
$O(\sqrt{\mathrm{cond}\,M})$ steps); at cavity scale use the shift-invert form above, or
better, (b).  For $\nu$ the code runs the power iteration on $M^{-1}A$ with an LU of $M$
--- the same once-per-mesh exception.  \emph{(b) No eigensolve at all.}  Only a
\emph{lower} bound on $\lambda_{\min}(M)$ is needed (a smaller value only pushes $Y$
outward), and finite elements hand one over for free: with $M=\sum_e P_e^{\!\top}M_eP_e$
the assembly of element mass matrices,
$x\hp Mx=\sum_e x_e\hp M_ex_e\ge\sum_e\lambda_{\min}(M_e)\|x_e\|^2
\ge\bigl(\min_e\lambda_{\min}(M_e)\bigr)\|x\|^2$, since every degree of freedom belongs to
at least one element; hence $\lambda_{\min}(M)\ge\min_e\lambda_{\min}(M_e)$ --- a loop over
elements with $3\times3$ (P2, 1-D) or $6\times6$ (P2, 2-D triangles) eigenproblems, no
global solve, and certified.  Likewise an \emph{upper} bound on $\nu$ suffices:
$\nu\le\|A\|_2/\lambda_{\min}(M)\le\sqrt{\|A\|_1\|A\|_\infty}/\lambda_{\min}(M)$, two
column/row sums --- pessimistic (it may push $Y$ out by an order of magnitude), but $Y$
only lengthens the quiet part of the tracked segment, which is cheap.  \emph{(c) Lumped
mass}: the smallest diagonal entry, exactly.
\paragraph{Singular mass matrix: a real gap in the derivation above.}
The bound presumes $M$ invertible, and the cavity's is not.  The discretised incompressible Navier--Stokes operator is a
differential-algebraic pencil: $M=\mathrm{diag}(M_v,0)$ (mass on velocity, nothing on
pressure), index 2.  Three consequences.  (i)~$\lambda_{\min}(M)=0$, $M^{-1}$ does not
exist, and the bound as written is void --- it must be re-derived, not patched.
(ii)~The pencil has infinite eigenvalues (as many as the algebraic constraints), which are
not eigenvalues of anything physical; $\det(sM-A)$ is a polynomial of degree below $n$,
the argument principle for $f$ still holds --- $f$ is a rational function whose finite
poles are the finite plant and controller eigenvalues --- but the tail closure is exactly
the part that needs the re-derivation.  (iii)~The repair: the physics lives on the
discretely divergence-free subspace.  Writing the DAE as
$M_v\dot v=A_{vv}v+B^{\!\top}p+b\,u$, $Bv=0$, and eliminating $p$ with the discrete Leray
projector $\Pi$ (the $M_v$-orthogonal projector onto $\ker B$; applying $\Pi$ is one
saddle-point, Stokes-type, solve), the velocity dynamics are an ODE,
$\dot v=\Pi M_v^{-1}A_{vv}v+\Pi M_v^{-1}b\,u$, and for a channel that forces the momentum
equation and measures velocity --- which is our case: $b$ is a body force (a boundary
actuation is recast as one by lifting), $c$ a velocity functional --- $G(s)$ is strictly
proper and the Neumann argument goes through verbatim on the subspace, with
$\nu:=\|\Pi M_v^{-1}A_{vv}\Pi\|$ (power iteration; each application one $M_v$ solve plus
one projection; once per mesh) and $\lambda_{\min}(M_v)$ of the velocity mass matrix,
which is SPD, so (b) applies to it unchanged.  Two design constraints fall out of this
and are worth recording now.  If the sensor measured \emph{pressure} (a common cavity
choice), $G$ would be proper but not strictly proper --- pressure answers a force
instantaneously --- so $G(\infty)\ne0$ and $f(s)\to1-KG(\infty)$ rather than $1$; the
winding argument survives provided the tail confines $f$ to a disk not containing the
origin, i.e.\ $|KG(\infty)|<1$ or a strictly proper controller ($K=0$).  And if an input
entered the \emph{constraint} equation, $G$ would carry a polynomial part and no tail
closure would exist --- so inputs must stay in the momentum equation, which is a
modelling rule, not a numerical one.  The projected-subspace version of the bound is
written up in \S\ref{sec:todo} as the first correction to make, and it should be
validated on a Stokes/lid-driven benchmark before the cavity.

\subsection{Extension to Navier--Stokes: singular mass, real arithmetic}
\label{sec:ns}


The extension to the cavity changes two things: the mass matrix becomes singular
(descriptor form), and the arithmetic becomes real.  Both are covered here; the
shift-invert-at-zero decision for $\lambda_{\min}(M_v)$ is recorded in \S\ref{sec:todo},
item~1.

\paragraph{Descriptor structure (recap of the repair).}  The cavity pencil is
$M=\mathrm{diag}(M_v,0)$, index-2: the tail bound is re-derived on the discretely
divergence-free subspace via the Leray projector $\Pi$, with
$\nu=\|\Pi M_v^{-1}A_{vv}\Pi\|$ and $\lambda_{\min}(M_v)$; the modelling rules are
\emph{inputs in the momentum equation only} (else $G$ grows a polynomial part and no
tail closure exists) and \emph{velocity outputs preferred} (pressure sensors leave $G$
proper but not strictly proper, shifting the tail disk to $1-KG(\infty)$, admissible only
if $|KG(\infty)|<1$ or $K=0$).  All once-per-mesh constants are computed on $M_v$, which
is SPD, so every recipe of this section applies unchanged.

\paragraph{Real arithmetic (what changes, all of it in our favour).}  The GL plant is
complex; the Navier--Stokes plant is \emph{real}, and realness is structure:
(i)~\emph{Conjugate symmetry of the spectrum.}  Eigenvalues come in pairs
$\lambda,\bar\lambda$; the rightmost eigenvalue is either real or a conjugate pair with
\emph{equal} real parts --- so the abscissa ``tie'' created by a complex pair is not a
kink at all (one constraint $\mathrm{Re}\,\lambda\le-\eps$ covers both members, since
$\mathrm{Re}\,\lambda=\mathrm{Re}\,\bar\lambda$ identically), and per-eigenvalue tracking
lists only the $\mathrm{Im}\ge0$ representatives.  (ii)~\emph{Frequency-response
symmetry.}  $\sigma_1(-\omega)=\sigma_1(\omega)$, so the band, the working set, the
certification sweep and the peak polish all live on $\omega\ge0$: half the frequencies,
half the certification cost, for free.  (iii)~\emph{Real factorisations.}  Away from the
frequency axis's complex shifts, the once-per-mesh objects ($M_v$, $A\hp A$, $\sigma M-S$)
are real symmetric: real Cholesky/Lanczos, roughly a factor $2$--$4$ cheaper in memory
and flops than their complex counterparts.  The frequency-axis LUs of $R(i\omega)^{-1}$
remain complex (the shift is imaginary), as before.  (iv)~\emph{Winding symmetry.}
$f(\bar s)=\overline{f(s)}$, so the phase increments on the lower half-line mirror the
upper ones: track $[\ell,\ \ell+iY]$ only and double it, adding the (real-axis) start
point once --- the certificate cost halves, again for free.  None of these changes any
statement in this note; they change constants, all downward.  (Added to the to-do list's
item~1 validation plan: check both symmetries explicitly at the Stokes benchmark, since
they are also free \emph{consistency tests} of the implementation.)
The controller scalar obeys the same logic one size down: from \eqref{eq:ks},
$|k(s)|\le|K|+\|n\|\,\|B_c\|/(|s|-\|A_c\|)$ for $|s|>\|A_c\|$.  Both bounds decay like
$1/|s|$, so their product can be driven below any threshold.

\paragraph{Choosing the endpoint $Y$, and why the answer is a certified integer.}
Choose the threshold $\tfrac12$ and let $Y$ solve the scalar inequality
\[
\Bigl(\sup_{|s|\ge Y}|k(s)|\Bigr)\cdot
\frac{\|c\|\,\|b\|}{\lambda_{\min}(M)\,(Y-\nu)}\;\le\;\frac12
\]
(monotone in $Y$; one bisection, once per design).  Beyond the endpoints, and on the whole
closing arc, $|f(s)-1|=|k(s)G(s)|<\tfrac12$: the curve $f$ is confined to the disk of
radius $\tfrac12$ around $1$.  Confinement does \emph{not} mean $f$ is constant there ---
it keeps moving --- but a curve inside that disk stays in the half-plane
$\{\mathrm{Re}\,f>\tfrac12\}$, its argument is confined to $(-\pi/6,\pi/6)$, and a curve
whose argument lives in a sector of width $<2\pi$ \emph{cannot wind around the origin}.
So the arc's net phase contribution is not zero by magic; it is a number provably in
$(-\pi/3,\pi/3)$.  The total over the \emph{closed} contour is exactly $2\pi\times$ an
integer --- winding numbers of closed curves are integers by topology, not by numerics ---
so the implementation takes the tracked segment phase $\Delta\varphi$, and returns the
nearest integer to $\Delta\varphi/2\pi$: the rounding is \emph{certified}, because the
unaccounted arc contribution is provably less than a sixth of a turn, far below the half
turn at which nearest-integer rounding could err.  Yes: truly an integer, and the round
is part of the proof, not a numerical apology.





On the cost of the factorisation itself: factoring a sparse matrix does create fill, but how
much is a matter of ordering and dimension, not a catastrophe.  In 1-D, $M$ is banded and
the factor is banded with \emph{zero} fill --- $O(n)$.  In 2-D, with a nested-dissection
ordering, the Cholesky factor of an $n$-node mesh carries $O(n\log n)$ nonzeros and
$O(n^{3/2})$ flops --- the same class as the LUs of $R(s)^{-1}$ the pipeline performs by
the hundreds, so ``ruins sparsity'' is too strong: it \emph{taxes} it, at the standard
rate, and once per mesh.  That said, for $M$ specifically there is a cheaper road: the
mass matrix is spectrally well conditioned (condition number $O(1)$--$O(10^2)$,
mesh-quality dependent, not $n$-dependent), so conjugate gradients solve $Mx=b$ in a few
tens of \emph{matvecs} with no factor at all --- and CG inside Lanczos-on-$M^{-1}$ is a
perfectly good once-per-mesh scheme.  So the efficiency ladder is: lumped diagonal
(free) $\to$ element bound (free, certified) $\to$ CG-based shift-invert (matvecs only)
$\to$ Cholesky-based shift-invert (fastest per iteration, pays fill once).
The element-wise bound, for its part, is unconditionally a lower bound.  The two inequalities in its
derivation are exact: $x_e\hp M_ex_e\ge\lambda_{\min}(M_e)\|x_e\|^2$ is the definition of
a smallest eigenvalue, and $\sum_e\|x_e\|^2\ge\|x\|^2$ holds because every global degree
of freedom appears in at least one element (shared dofs are counted \emph{more} than
once, which only helps).  No mesh assumption, no asymptotics --- a certified floor.

Two decisions close this topic.  \emph{$\lambda_{\min}(M)$}: shift-invert
(Cholesky- or CG-based) for the tight value, with the element bound retained as the
certified cross-check --- a tight constant shortens the tracked segment on every future
design, so the one-time cost repays itself immediately.  \emph{$\|A\|_2$}: yes ---
Lanczos on the Hermitian $A\hp A$ (matvecs only, extremal target, a few dozen iterations)
gives the true $\|A\|_2=\sqrt{\mu_{\max}}$, tighter than
$\sqrt{\|A\|_1\|A\|_\infty}$ by typically a small factor but occasionally much more on
convection-dominated operators; then $\nu\le\|A\|_2/\lambda_{\min}(M)$, or, tightest of
all, the direct power iteration on $M^{-1}A$ already described.  The recipe recorded in
\S\ref{sec:todo} (item~2, updated): tight values by shift-invert/Lanczos, certified
bounds (element-wise; H\"older) computed alongside as floors/ceilings --- the pipeline
uses the tight value and \emph{checks} it against the bound, so a bug in either is caught
by the other.

$P_{\mathrm{ctrl}}$ counts controller poles right of the line:
$P_{\mathrm{ctrl}}(\ell)=\#\{\lambda\in\mathrm{eig}(A_c):
\mathrm{Re}\,\lambda>\ell\}$, an exact $r\times r$ computation.  One nuance worth
noticing: the production line is $\ell=-\eps/2<0$, so a \emph{stable but slow} controller
pole with $\mathrm{Re}\,\lambda\in(-\eps/2,0)$ is counted in $P_{\mathrm{ctrl}}$ --- and
correspondingly contributes a pole of $f$ inside the region --- even though nothing is
unstable.  The count formula balances this automatically; the point is only that
``poles right of the line'' is not a synonym for ``unstable poles''.
where $P_{\mathrm{plant}}$ (plant eigenvalues right of the line) is
\emph{design-independent} --- computed once per mesh, and at cavity scale supplied from

\paragraph{Obtaining $P_{\mathrm{plant}}$ certified at cavity scale.}
A one-time cost per mesh and Reynolds number, amortised over the entire campaign; the
recipe has three steps.  (1)~\emph{Confine}: bound the region
where right-of-line eigenvalues can live.  The numerical range gives it cheaply:
$\mathrm{Re}\,\lambda\le\max_x \mathrm{Re}(x\hp Ax)/(x\hp Mx)$, a Hermitian eigenvalue
problem for the symmetric part (one Lanczos run), and similarly for $|\mathrm{Im}\,\lambda|$
--- together a certified bounded rectangle to the right of the line.  (2)~\emph{Count
inside the rectangle}: a contour-integral eigenvalue count (Beyn/FEAST style: the number
of eigenvalues inside a contour equals the rank of a contour integral of the resolvent,
estimated with a block of random probes and quadrature on the rectangle's boundary ---
every quadrature node is one plant factorisation, reused for all probes), cross-validated
at two quadrature orders and two probe blocks.  (3)~\emph{Polish}: each counted eigenvalue
individually by shift-invert to machine precision, so the count is backed by exhibited
eigenvalues.  For a cavity at moderately supercritical conditions the physical count is
small and known from the global-instability literature, so step~(2) is a confirmation
priced at $O(10^2)$ factorisations --- once, not per design.

Step~(1) deserves care, because it is easy to misread as the non-Hermitian pencil
$Ax=\lambda Mx$ --- it is not.  It is a Hermitian problem.
Write $S=\tfrac12(A+A\hp)$ for the Hermitian part; if $Ax=\lambda Mx$ then
$\lambda=x\hp Ax/x\hp Mx$ and
$\mathrm{Re}\,\lambda=x\hp Sx/x\hp Mx\le\mu_{\max}$, where $\mu_{\max}$ is the largest
eigenvalue of the Hermitian pencil $Sx=\mu Mx$ (the numerical range, or field of values,
of $A$ in the $M$ inner product).  Hermitian pencils are where Lanczos is at home, in
three respects.  \emph{Finding the top}: extremal eigenvalues converge first; run
regular-mode Lanczos on $M^{-1}S$ (the once-per-mesh Cholesky of $M$ again), or ---
faster when the top of the spectrum is clustered, which for a convection--diffusion
operator it is, the top being the least-damped end --- shift-invert at a shift $\sigma$
to the \emph{right} of the spectrum, where $\sigma M-S$ is positive definite, hence
Cholesky-factorisable, and Lanczos on $(\sigma M-S)^{-1}M$ returns the eigenvalues
nearest $\sigma$, i.e.\ the top ones, in a handful of steps.  \emph{Certifying the top}:
Rayleigh--Ritz gives a Ritz value $\theta\le\mu_{\max}$ --- a lower bound, the wrong side
for confinement --- but a certified \emph{upper} bound costs one factorisation: pick
$\sigma=\theta+\delta$ and attempt the sparse $LDL\hp$ of $\sigma M-S$; if every pivot is
positive, Sylvester's inertia theorem says $\sigma M-S\succ0$, hence $\mu_{\max}<\sigma$.
\emph{A Cholesky that succeeds is a proof.}  The same two moves on the skew part
$\tfrac1{2i}(A-A\hp)$ bound $|\mathrm{Im}\,\lambda|$, and the rectangle is certified.
Singular $M$: on the velocity subspace with $\Pi$, as in the previous answer.
So step~(1) is a rectangle with a certificate, and it answers the one question no mesh of
shifts can --- \emph{how far right and how far up must we look?}  The common practice of
shift-invert on a mesh of imaginary-axis shifts is step~(3) of this recipe, and it does
yield the count \emph{if} it finds every eigenvalue; what it cannot supply is the word
``every''.  Shift-invert Arnoldi at
$\sigma$ returns the eigenvalues nearest $\sigma$, not a radius within which it has seen
everything; two neighbouring shifts leave a gap about which the method is silent, and an
eigenvalue far from any shift (large real part, or large imaginary part beyond the last
shift) is simply invisible.  Step~(2) is what turns ``we found $K$'' into ``there are
$K$''.  On the fluid-mechanics habit: it is a reasonable one, because the physical modes
are few, isolated and known in advance --- and that is exactly the regime in which
step~(2) is a cheap confirmation rather than a search.  There is also a safety net
specific to our use of the number: $P_{\mathrm{plant}}$ is a constant added to every
design's count, so an error in it is a constant offset in every certificate --- detected
by the first design whose stability is known by other means (the continuation gate, or at
benchmark scale the dense eigensolve): a design known to be stable would read a nonzero
count.  Honest status: at GL scale the shipped \code{plant\_pole\_data} is dense (and
refuses $n>3000$, asking for $P_{\mathrm{plant}}$ and the near-line modes explicitly);
steps (1)--(3) are the cavity recipe, written but not yet run.

The contour-counting idea has its own literature, added to the reference list:
T.~Sakurai \& H.~Sugiura, ``A projection method for generalized
eigenvalue problems using numerical integration'', \emph{J.\ Comput.\ Appl.\ Math.}\ 159
(2003) 119--128; E.~Polizzi, ``Density-matrix-based algorithm for solving eigenvalue
problems'', \emph{Phys.\ Rev.\ B} 79 (2009) 115112 (the FEAST solver); W.-J.~Beyn, ``An
integral method for solving nonlinear eigenvalue problems'', \emph{Linear Algebra Appl.}\
436 (2012) 3839--3863.  Yes, it is the argument principle, in matrix form.  Scalar
version: $\tfrac1{2\pi i}\oint\varphi'/\varphi\,dz=Z-P$ for a meromorphic $\varphi$.
Matrix version: the resolvent $R(z)=(zM-A)^{-1}$ is meromorphic with simple poles at the
eigenvalues (assumed simple) and residues $v_jw_j\hp/(w_j\hp Mv_j)$, so by the residue
theorem
\[
\frac1{2\pi i}\oint_\Gamma (zM-A)^{-1}M\,dz
\;=\;\sum_{\lambda_j\in\mathrm{int}\,\Gamma}\frac{v_jw_j\hp M}{w_j\hp Mv_j}
\;=\;P_\Gamma ,
\]
the spectral projector onto the eigenspace of the eigenvalues inside $\Gamma$, whose rank
--- or trace --- is the count.  And the trace is
$\tfrac1{2\pi i}\oint\mathrm{tr}[(zM-A)^{-1}M]\,dz
=\tfrac1{2\pi i}\oint\tfrac{d}{dz}\log\det(zM-A)\,dz$: literally the argument principle
for $\varphi=\det(zM-A)$.  That answers ``what about zeros'': for the determinant the
zeros \emph{are} the eigenvalues (with multiplicity) and there are no poles at all ---
$\det(zM-A)$ is a polynomial --- so Beyn/FEAST counts zeros only, nothing to subtract.
Contrast with $f$ in \eqref{eq:count}: it has poles as well, because we divided the plant
and controller determinants out to reach a scalar; that division is what buys a
one-dimensional contour at the cost of two pole counts.  For $P_{\mathrm{plant}}$ no
such low-rank reduction exists, hence the matrix integral, computed with probes: for a
random $Y\in\C^{n\times m}$ with $m$ at least the expected count,
$A_0=\tfrac1{2\pi i}\oint R(z)MY\,dz$ has rank equal to the count (generically), and
$A_1=\tfrac1{2\pi i}\oint zR(z)MY\,dz$ then yields the eigenvalues themselves from the
$m\times m$ pencil $(A_1,A_0)$ restricted to the range --- Beyn's method, so steps (2)
and (3) can share the factorisations.  Quadrature: Gauss or trapezoid on each edge of the
rectangle; the error decays exponentially in the number of nodes for eigenvalues at a
distance from the contour and slowly for eigenvalues near it --- which is why the count
is cross-validated at two orders, and why the rectangle's left edge should sit at
$\ell-\delta$ rather than on the line, letting the polish of step (3) decide the side of
$\ell$ for anything found near it.  Rigour, honestly: the rank of $A_0$ is decided from a
singular-value gap, a numerical judgement, not a theorem; the standard of rigour is
``count $K$ at two quadrature orders and two probe blocks, then exhibit $K$ eigenvalues
inside with residuals at machine precision'' --- a certificate by exhibition and
consistency, of the same kind as the disk sweep of \S\ref{sec:cert}.
\subsection{Contour counting from scratch: the argument principle in matrix form}
\label{sec:contour}


The construction, in five steps, assuming nothing.

\emph{Step 1: one eigenvalue, one integral.}  Let $\lambda_0$ be a simple eigenvalue of
$Ax=\lambda Mx$ with right/left eigenvectors $v,w$.  Near $\lambda_0$ the resolvent
behaves like a simple pole with a rank-one residue:
$R(z)=(zM-A)^{-1}=\dfrac{v\,w\hp}{(z-\lambda_0)\,w\hp Mv}+(\text{analytic near }\lambda_0)$.
(To check the residue: multiply by $(zM-A)$ and use $Av=\lambda_0Mv$, $w\hp A=\lambda_0w\hp M$.)
So integrating around a small circle $\Gamma_0$ enclosing only $\lambda_0$, the analytic
part contributes nothing and the residue theorem leaves
$\tfrac1{2\pi i}\oint_{\Gamma_0}R(z)M\,dz=\dfrac{v\,w\hp M}{w\hp Mv}$ --- a rank-one
matrix that maps any vector onto the eigendirection $v$.

\emph{Step 2: many eigenvalues.}  For a contour $\Gamma$ enclosing eigenvalues
$\lambda_1,\dots,\lambda_K$, residues add:
\[
P_\Gamma:=\frac1{2\pi i}\oint_\Gamma R(z)M\,dz
=\sum_{j\le K}\frac{v_jw_j\hp M}{w_j\hp Mv_j}.
\]
This $P_\Gamma$ is a \emph{projector} ($P_\Gamma^2=P_\Gamma$: multiply two terms and use
the biorthogonality $w_i\hp Mv_j=0$ for $i\ne j$): it projects onto the span of the
enclosed eigenvectors.

\emph{Step 3: why trace $=$ rank $=$ count.}  A projector's eigenvalues are only $0$ and
$1$ (any eigenvalue $\pi$ of $P$ obeys $\pi^2=\pi$).  Its \emph{rank} is the dimension of
the space it projects onto --- the number of enclosed eigendirections, i.e.\ the count
$K$.  Its \emph{trace} is the sum of its eigenvalues --- $K$ ones plus zeros $=K$ again.
Trace and rank are two roads to the same integer; the trace is the analyst's road, the
rank the algorithm's (step 5).

\emph{Step 4: the trace is the scalar argument principle in disguise.}  Jacobi's formula
--- the derivative of a determinant --- says
$\tfrac{d}{dz}\det X(z)=\det X(z)\cdot\mathrm{tr}\bigl[X(z)^{-1}X'(z)\bigr]$; proof in
two lines:
$\det(X+\epsilon X')=\det X\,\det(I+\epsilon X^{-1}X')
=\det X\,\bigl(1+\epsilon\,\mathrm{tr}[X^{-1}X']+O(\epsilon^2)\bigr)$, because
$\det(I+\epsilon B)=\prod_i(1+\epsilon\beta_i)=1+\epsilon\sum_i\beta_i+O(\epsilon^2)$.
With $X=zM-A$, $X'=M$:
$\mathrm{tr}[R(z)M]=\tfrac{d}{dz}\log\det(zM-A)$ --- and it
identifies $\mathrm{tr}\,P_\Gamma$ with $\tfrac1{2\pi i}\oint\varphi'/\varphi\,dz$ for
$\varphi=\det(zM-A)$: the scalar argument principle, counting the polynomial's zeros
(the eigenvalues, with multiplicity), no poles to subtract.

\emph{Step 5: computing without the trace --- probes.}  Forming $\mathrm{tr}[R(z)M]$
exactly would need $n$ solves per quadrature node.  Instead hit the projector with a thin
random block $Y_{\!p}\in\C^{n\times m}$, $m\gtrsim K$:
$A_0=P_\Gamma Y_{\!p}$ (per node: one factorisation, $m$ solves).  Since $P_\Gamma$ has
rank $K$, generically $A_0$ has rank exactly $K$, read off from its singular values (a
clean gap after the $K$-th).  That is the \emph{count}.  For the eigenvalues themselves,
one more moment, $A_1=\tfrac1{2\pi i}\oint zR(z)MY_{\!p}\,dz$, and the tiny $m\times m$
pencil built from $(A_1,A_0)$ has the enclosed $\lambda_j$ as its eigenvalues --- Beyn's
method.

\emph{Beyn or FEAST?}  Same integral, different engines.  \textsc{Beyn} (2012) is the
\emph{direct} method just described: two moments, one SVD, one tiny pencil ---
non-iterative, ideal when $K$ is small and one wants count-plus-eigenvalues in one shot.
\textsc{FEAST} (Polizzi 2009 --- the proper name of his solver and library, not an
acronym) uses the same contour integral as a \emph{filter} inside subspace iteration:
apply $P_\Gamma$ repeatedly to refine a subspace, Rayleigh--Ritz inside --- iterative,
more robust to quadrature error and to underestimating $m$, preferred when many interior
eigenvalues are wanted.  For $P_{\mathrm{plant}}$ --- few eigenvalues, one shot, count
first --- \textbf{Beyn is the right tool}, and the to-do item now says so
(\S\ref{sec:todo}, item~3).


\paragraph{Placing the contour.}
\emph{$\ell$} is the counting line's real part --- the same $\ell$ as in
\eqref{eq:count}; production value $\ell=-\eps/2$.  \emph{$\delta$} is the standoff of
the rectangle's \emph{left} edge from that line: the edge sits at $\ell-\delta$ so that
near-line eigenvalues fall strictly \emph{inside} the contour rather than on it, and the
shift-invert polish decides on which side of $\ell$ each of them lies.  \emph{The
contour}: right and top edges from the certified numerical-range rectangle (inertia
certificate) plus a margin; bottom edge its mirror (or the real axis for a real plant,
\S\ref{sec:ns}); left edge at $\ell-\delta$.  \emph{$Y$ and $\nu$, unified}: yes, the
same family --- $\nu$ confines the whole spectrum to the disk $|\lambda|\le\nu$ (a
circumscribed-disk statement), while the numerical-range rectangle is the sharpened,
per-direction version of the same confinement; the rectangle's $|\mathrm{Im}|$ bound can
therefore \emph{tighten} the winding endpoint $Y$ too, and the cleanest implementation
computes the rectangle once and derives both the contour and $Y$ from it, keeping the
Neumann-$\nu$ formula as the fallback when only the crude constants exist.
\emph{Quadrature robustness}: the danger: an eigenvalue (or a
region of large resolvent norm, the pseudospectrum of a non-normal $A$) close to an edge
slows the quadrature's convergence.  Three defences, layered: (i)~\emph{certified
standoff}: sweep each contour edge with the eigenvalue-free disks of \S\ref{sec:cert}
--- a 1-D sweep per edge, the same tool as the frequency band --- which certifies a
minimum spectral distance from the contour and, if a disk refuses to certify, moves the
edge before any quadrature is attempted; (ii)~\emph{resolvent monitoring}: the quadrature
already evaluates $R(z)MY_{\!p}$ at every node --- its norm is read off for free, and a
spike flags a near-pseudospectral passage even when eigenvalues are distant
(non-normality), triggering edge displacement or local refinement; (iii)~the
\emph{two-order cross-validation} of the count, which is the net under everything.
Extra margin: yes, cheap and wise --- the rectangle is enlarged by a fixed fraction
before quadrature, since a larger contour costs nothing (the count inside is unchanged)
while a tight one courts exactly that slow-convergence risk.


\paragraph{Why a mesh of shifts is not a certificate --- and how to make it one.}
Shift-invert on a mesh of imaginary-axis shifts (perhaps with a second line at
$\mathrm{Re}\,s=1$) is what everyone does, and what the shipped code effectively does at
GL scale (with a dense solve standing in for the mesh); it is not a certificate --- for
one reason that step~(1) fixes and one that step~(2) fixes.  Without (1) the mesh has no
\emph{edges}: nothing establishes that the shifts need not go beyond $|\mathrm{Im}\,s|=60$
or $\mathrm{Re}\,s=1$, and a spurious or physical mode at $(2,200)$ is invisible to both
lines.  The numerical-range rectangle is the only cheap object that answers ``how far'',
and it is one Hermitian eigensolve plus one Cholesky.  Without (2) the mesh has no
\emph{radii}: shift-invert at $\sigma$ sees the nearest eigenvalues, and says nothing
about the strip between two neighbouring shifts, so a second line at $\mathrm{Re}\,s=1$
does not confirm anything --- it adds a second row of blind spots, at half the distance.
There is, however, a way to make exactly that scheme rigorous, with a tool we already
own: the eigenvalue-free disks of \S\ref{sec:cert}.  At shift $\sigma$, the
distance $r(\sigma)$ to the nearest eigenvalue certifies the open disk of radius
$r(\sigma)$ empty (with the same hedge as in \S\ref{sec:cert}: request a few
eigenvalues, so that Arnoldi's ``nearest'' is trustworthy); cover the rectangle by such
disks and the count is the number of eigenvalues exhibited.  That is the Schwerdtner
sweep run over an \emph{area} instead of a line, and its cost is the catch: the number of
disks grows like $(\text{area})/r^2$, and near the line the radii are tiny (the leading
GL mode sits $0.007$ from the production line), so the hot regions eat factorisations ---
whereas the contour count of (2) costs $O(10^2)$ factorisations regardless of where the
eigenvalues sit.  Verdict: keep (1), because nothing else bounds the search; take (2) as
either the contour count or the disk covering, whichever one trusts more --- both are the
argument principle, one on the boundary and one by exhaustion of the interior; and use
the mesh, and the second line, as what they are good at: shifts for the polish of (3),
the second line being a sensible place to polish anything the count locates far to the
right.  And the fluid-mechanics practice of trusting the mesh is, seen this way, a
covering argument supplied by physical prior knowledge --- legitimate for a plant one
knows, which is exactly why $P_{\mathrm{plant}}$, computed once and cross-checked once,
is the least worrying constant in the whole certificate.

\paragraph{How many eigenvalues per shift ($\code{nev}$).}
The campaign measured $\code{nev}=1$ adequate \emph{at GL scale under the
validation harness}: every disk cross-checked against dense ground truth, $104/104$ ---
so ``$\code{nev}=1$ is enough here'' is a measured fact, and paying $3\times$ per shift
for hypothetical robustness was rightly refused.  The opposite hedge --- request a few --- addresses the regime without ground truth
(cavity), where a misconverged ``nearest'' silently inflates a disk.
The shipped-and-recommended policy: $\code{nev}=1$ by default, with \emph{escalation on
suspicion} --- re-solve with $\code{nev}=3$ only when the returned pair's residual is
loose, when two adjacent disks disagree about their overlap, or when the disk is
critical (its radius decides a certificate rather than merely advancing the sweep).
Cheap where it can be, careful where it must be; at GL the escalation almost never
fires, which is itself the measurement that justifies the default.
the known physical instability structure --- and $P_{\mathrm{ctrl}}$ is an exact $r\times
r$ computation.  Each sample of $f$ along the line costs one plant LU
($G=c\hp Rb$); $f'$ costs one extra solve on the same factorisation
($G'=-c\hp R\,M\,R\,b$), enabling phase-tracking with derivative information (its use in
practice: the question below).  The arc at infinity is closed analytically by the tail
bound derived above.  
\paragraph{Using the derivative information.}
Two uses.  \emph{Predictor}: the local phase speed is
$\mathrm{d}\arg f/\mathrm{d}s=\mathrm{Im}(f'/f)$, so the next sample is placed at
$\Delta s\approx\varphi_{\mathrm{budget}}/|f'/f|$ with a budget of a fraction of $\pi$ ---
adaptive stepping by prediction instead of by trial-and-refine.  \emph{Validator}: the
tracked increment is cross-checked against the Hermite estimate
$\tfrac{\Delta s}2[\mathrm{Im}(f'/f)_{\mathrm{left}}+\mathrm{Im}(f'/f)_{\mathrm{right}}]$;
disagreement triggers refinement before the midpoint test even runs.  Cost: $f'$ is one
extra triangular solve on the \emph{same} factorisation --- at GL scale that is a
$\sim30\%$ overhead per sample, at cavity scale (where the LU dominates the solve by an
order of magnitude) nearly free --- against which the predictor cuts the sample count by
a factor that was worth it at GL already and is decisive wherever each sample carries a
fresh factorisation.  So: worth it, increasingly so with scale.
Sampling is adaptive with three hardening
mechanisms (mandatory resolution near the poles of $f$, midpoint validation of accepted
steps, and a mandatory knot grid), each added after a measured failure; the certificate
reproduces dense eigenvalue counts exactly on $104/104$ validation designs at

\paragraph{Cost at cavity scale, honestly.}
At cavity scale one factorisation is the dominant unit
(minutes, gigabytes), so $500$--$900$ of them serially is hours --- and
three facts make it acceptable and a fourth makes it smaller.  (1)~Cadence: the winding
certificate runs per \emph{saved} design --- a handful per case --- never per step.
(2)~Parallelism: samples are independent; $64$ workers make it $\sim10$ rounds, and the
memory footprint is one factorisation \emph{per worker at a time} (each worker processes
its queue serially and discards), not one per sample --- the footprint objection
dissolves.  (3)~The derivative-informed predictor above cuts the sample count.
(4)~Structural reduction, if still needed: the samples concentrate near the poles of $f$;
away from the hot regions the same Neumann logic that closes the arc certifies long
quiet stretches of the \emph{segment} from a few samples with interval bounds on
$|kG|$ --- a coarse-certify/fine-track split that brings the count toward $O(10^2)$.
Projection, to be validated at GL first (to-do list of the optimisation note): tens of
minutes per saved design on a modest allocation --- a final-answer cost, not a loop
cost.




$\approx500$--$900$ plant solves per design.  The production line is $\ell=-\eps/2$: a
certified count of zero there is an \emph{unconditional} proof of margin-stability, with
no continuation premise.

%======================================================================
\section{Gradient of the abscissa --- and do we need its subgradients?}
\label{sec:subgrad}

\subsection{The gradient}
Let $\lambda$ be the rightmost eigenvalue, assumed \emph{simple}, with right and left
generalized eigenvectors
$\Acl v=\lambda \Mcl v$ and $w\hp\Acl=\lambda\,w\hp\Mcl$.  Standard eigenvalue
perturbation gives
\begin{equation}
\boxed{\;
\delta\lambda \;=\; \frac{w\hp\,(\delta\Acl)\,v}{\,w\hp \Mcl\, v\,},
\qquad
\delta\alpha=\mathrm{Re}\,\delta\lambda ,
\;}
\label{eq:alphagrad}
\end{equation}
which has \emph{exactly} the shape of the master formula \eqref{eq:master} with
$(\tilde u,\tilde v)$ replaced by $(w/\overline{(w\hp\Mcl v)},\,v)$: the same
group-by-group table of \S\ref{sec:grads} applies verbatim, so the full gradient of
$\alpha$ with respect to every design parameter costs \emph{one eigenpair} --- obtained by
shift-invert iteration and its adjoint on one plant factorisation.  (The denominator
$w\hp\Mcl v$ also serves as a health monitor: it tends to zero when the eigenvalue
approaches a defective (Jordan) configuration, where $\alpha$ ceases to be Lipschitz;
non-generic, but worth watching.)

\subsection{Do we need subgradients?}
Three answers, in increasing order of care.

\emph{(i) Generically, no.}  For a complex (non-real) system like ours there is no
automatic conjugate pairing, so the rightmost eigenvalue is generically unique and simple;
$\alpha$ is then differentiable and \eqref{eq:alphagrad} is its honest gradient.  This was
the measured situation even in the hardest case of the campaign: the design pinned on the
margin (K+B+C) carried a \emph{single} eigenvalue at the line --- we tested the
multiple-eigenvalue hypothesis explicitly and rejected it.

\emph{(ii) At ties, the subdifferential exists but we avoid needing it.}
$\alpha=\max_j\mathrm{Re}\,\lambda_j$ is a max, so when several eigenvalues tie at the
rightmost real part the correct first-order object is the convex hull of their individual
gradients \eqref{eq:alphagrad} --- structurally identical to the frequency-max story of
the optimisation note.  And the remedy is the same for the third time: \emph{never take
the max; name the members}.  The diagnostic KKT tool already treats each near-line
eigenvalue as its own function (gradients by matched perturbation of the tracked
eigenvalues, multipliers by nonnegative least squares); the planned explicit-constraint
inner does likewise, imposing $\mathrm{Re}\,\lambda_j(\theta)\le-\eps$ \emph{per tracked
eigenvalue}.  Each such constraint is smooth with gradient \eqref{eq:alphagrad}; the tie
lives in the active set, where the solver is competent, and the subdifferential of the max
is never formed.

\emph{(iii) The genuinely nonsmooth points are defective eigenvalues,} where two
eigenvalues collide into a Jordan block and $\alpha$ behaves like a square root
(non-Lipschitz; no useful subgradient).  These are avoided rather than modelled: the
margin formulation keeps the optimiser $\eps$ away from the axis where the physical
coalescences live, and $|w\hp\Mcl v|$ flags any approach.

The parallel with the singular-value story is close but not exact, and the difference is
worth owning.  The exact analogue of the
$\sigma_1=\sigma_2$ crossing is case~(ii) above --- a \emph{semisimple} tie, a max of two
smooth branches, a Lipschitz kink, handled by enrollment.  Case~(iii) --- the defective
collision --- is strictly \emph{worse} than anything singular values can do, and here is
the pleasant reason why: the $\sigma$'s are eigenvalues of the \emph{Hermitian} operator
$T\hp T$, and Hermitian eigenvalues are always semisimple --- crossings, yes; Jordan
blocks, never.  Eigenvalues of the non-normal $\Acl$ enjoy no such protection, which is
exactly why the abscissa needs the extra care of case~(iii) while the frequency response
never does.  Non-normality, the physical heart of this whole flow problem, is also the
algebraic reason its stability constraint is the more delicate object.

In short: \textbf{value + gradient of individually tracked eigenvalues suffice}; the max
structure is handled by enrollment, exactly as $\sigma_2$-enrollment handles
singular-value ties, and no subdifferential machinery is required anywhere in the design.

%======================================================================
\section{How \S5--\S\ref{sec:cert} compare with Benner--Mitchell (2018)}
\label{sec:bm}


The comparison is natural because their paper and our oracle share an architecture.  Their starting point is the classical level-set method of
Boyd--Balakrishnan and Bruinsma--Steinbuch (``BBBS'', 1990): bisect on the level
$\gamma$, and at each level pay one \emph{full Hamiltonian eigendecomposition} to find
all frequencies where $\sigma_1=\gamma$ --- robust, global, and $O(n^3)$ \emph{per
level}.  Their observation: that work distribution is upside-down --- cheap local
information is being bought with the most expensive global tool.  Their method
rebalances: run standard \emph{optimization} in $\omega$ to convergence on local peaks
--- this is where a reader may recall ``splines'': the local maximisation is
accelerated by cubic/Hermite \emph{interpolation} of $\sigma_1(\omega)$ using values
and derivatives (a refinement lineage their paper traces back through the level-set
literature), giving superlinear local convergence for a few evaluations --- and invoke
the expensive global eigendecomposition only to \emph{verify} the converged candidate
(typically once or twice per norm computation instead of once per bisection level).
Faster, and, as their title says, more accurate: the optimizer polishes to full
precision, which plain level-set bisection does not.
Now place \S5--\S\ref{sec:cert} beside it.  The division of labour is
\emph{identical} --- optimise locally, certify globally, and arrange for the certificate
to be consulted as rarely as possible: our Newton polish in $\omega$ (\S5) is their
interpolation-accelerated local stage (same derivative information, same superlinear
character), and our certified disk sweep (\S\ref{sec:cert}) plays their verification

Our local stage differs from theirs in one respect: we use Newton on
$\partial_\omega\sigma_1$.  The two are close relatives: Hermite interpolation through
$(\sigma_1,\partial_\omega\sigma_1)$ at two bracketing points is a secant-flavoured
Newton that borrows information across the bracket, while our Newton uses one point and a
second derivative.  Asymptotically Newton is faster (quadratic versus superlinear); far
from the peak the interpolation is more robust (it cannot propose a step outside its
bracket).  The choice is dictated by the starting point each method actually has: they
start cold, from level-set midpoints, so robustness wins; our peaks are warm-started from
the previous outer round, where the peak has barely moved, so Newton sits inside its
basin from the first iterate and the quadratic rate is real.  Both spend $2$--$4$ oracle
calls per peak; the difference is marginal, the contexts decide, and a safeguarded hybrid
(Newton with a Hermite-bracket fallback when $\partial^2_\omega\sigma_1\ge0$ signals a
plateau) is the best of both --- noted for the next iteration.
stage.  The two differences are deliberate.  \emph{First, the global tool itself}: their
verifier is the dense Hamiltonian eigendecomposition --- exact, but $O(n^3)$ and
dense-only; ours is the eigenvalue-free-disk certificate of Schwerdtner--Mengi--Voigt,
which needs only sparse shift-invert solves, i.e.\ (P1)--(P3), and therefore survives
the trip to cavity scale.  So: no, they do not use the Schwerdtner-style certification
--- at their scales they do not need it; we adopt it \emph{because} their verifier is
the one component of the architecture that does not scale.  \emph{Second, the
amortisation horizon}: they optimise and verify within \emph{one} norm computation; our
exchange loop amortises across an entire \emph{design} trajectory --- the working set
$\Omega$ carries candidate peaks from design to design, warm-starting both the local
stage (in $\omega$ \emph{and} in $\theta$) and the bisection of the certificate from the
incumbent level.  In one sentence: \emph{our oracle follows the Benner--Mitchell division of labour ---
cheap local optimisation, rare global verification --- with Newton in place of their
interpolation, their dense verifier replaced by the Schwerdtner--Mengi--Voigt
certificate, and the amortisation horizon stretched from one norm computation to the
whole optimisation}.

(``Per-norm'' means per computation of one design's $H_\infty$ norm; ``amortise'' means: spread the cost of the expensive component over many uses of its result.  They
amortise \emph{within} one norm computation --- one expensive global verification serves
many cheap local steps.  We keep that, and add a second horizon: \emph{across designs} ---
the working set $\Omega$, the singular vectors, the peak locations and the incumbent
level all carry over from one design to the next, so the $k$-th norm evaluation starts
where the $(k-1)$-th finished and costs a fraction of a cold one.  On ``architecture'':
fair objection --- the word claimed too much, since our local stage is Newton, not their
interpolation; the sentence above now says ``division of labour'' and names the
substitution explicitly.

\paragraph{References.}
Boyd \& Balakrishnan (1990); Bruinsma \& Steinbuch (1990) --- the BBBS level-set method.
Guglielmi, G\"urb\"uzbalaban \& Overton, SIMAX 2013 --- spectral-value-set iteration,
the fast local (non-certified) large-scale line.
Benner \& Voigt, \emph{Math.\ Control Signals Syst.}\ 2014 --- structured pseudospectral
$H_\infty$ computation for large-scale descriptor systems.
Aliyev, Benner, Mengi, Schwerdtner \& Voigt, SIMAX 2017 --- greedy subspace method for
large-scale $L_\infty$ norms.
Benner \& Mitchell, SISC 2018 --- the optimization-rebalanced norm computation above.
\revQ{P.~Schwerdtner, E.~Mengi \& M.~Voigt, ``Certifying global optimality for the
$L_\infty$-norm computation of large-scale descriptor systems'', \emph{IFAC-PapersOnLine}
53(2) (2020), 21st IFAC World Congress --- the eigenvalue-free-disk certificate of
\S\ref{sec:cert}; note that Mengi is a co-author of both this certificate and the
subspace methods of \S\ref{sec:proj}, which share the same view of the problem.}
Apkarian \& Noll, IEEE TAC 2006 (with the multiband and 2015 parametric papers, and
Apkarian's 2018 infinite-dimensional structured $H_\infty$) --- the synthesis line,
treated at length in the optimisation note.  Curtis, Mitchell \& Overton 2017 --- the
inner solver.

\paragraph{References for the contour counting.}
Sakurai \& Sugiura, \emph{J.\ Comput.\ Appl.\ Math.}\ 159 (2003) 119--128; Polizzi,
\emph{Phys.\ Rev.\ B} 79 (2009) 115112 (FEAST); Beyn, \emph{Linear Algebra Appl.}\ 436
(2012) 3839--3863 --- contour-integral eigenvalue counting and extraction, the matrix
form of the argument principle used for $P_{\mathrm{plant}}$ in \S\ref{sec:winding-num}.

\revP{\section{Projection: a reduced model for the inner loop}
\label{sec:proj}

Everything before this section evaluates the oracle by one sparse factorisation per
frequency, and the cost accounting of \S\ref{sec:cost} is built on that unit.  This
section describes a different organisation of the same algebra --- a low-dimensional
model of the resolvent that serves every frequency, and every design, from a handful of
factorisations --- states what has been measured on GL, and what remains to be measured
on the cavity.  It rests on a published method, which is described first.

\subsection{The method of Aliyev, Benner, Mengi, Schwerdtner and Voigt (2017)}
\label{sec:proj-paper}
Reference: N.~Aliyev, P.~Benner, E.~Mengi, P.~Schwerdtner, M.~Voigt, ``Large-scale
computation of $\mathcal L_\infty$-norms by a greedy subspace method'', \emph{SIAM J.\
Matrix Anal.\ Appl.}\ 38(4) (2017) 1496--1516 (arXiv:1705.10086).  The setting is more
general than a descriptor system: an $\mathcal L_\infty$-function
\[
H(s)=C(s)\,D(s)^{-1}B(s),\qquad
D(s)=\sum_{j=1}^{\kappa_D}h_j(s)D_j,\quad B(s)=\sum_jf_j(s)B_j,\quad C(s)=\sum_jg_j(s)C_j,
\]
with $D_j\in\C^{n\times n}$ large and sparse, $B_j$ tall-and-skinny, $C_j$ short-and-fat,
and $f_j,g_j,h_j$ meromorphic scalar functions --- descriptor and delay systems are the
paper's examples.  \emph{Our closed loop is natively of this form}: by
\S\ref{sec:solve}, $D(s)=s\,M-A-k(s)\,bc\hp$, i.e.\ $\kappa_D=3$ with
$h=(s,1,-k(s))$ and $D=(M,A,bc\hp)$, the controller entering only through the scalar
meromorphic coefficient $k(s)$ of \eqref{eq:ks}.  Note the $\mathcal L_\infty$ in the
title: the supremum of $\sigma_1$ over the imaginary axis is defined whether or not the
system is stable, and the method never assumes stability.

\emph{Reduction.}  Choose $V,W\in\C^{n\times\tilde n}$, $\tilde n\ll n$, and project the
middle factor two-sidedly (Petrov--Galerkin):
$\tilde H(s)=\tilde C(s)\tilde D(s)^{-1}\tilde B(s)$ with $\tilde D_j=W\hp D_jV$,
$\tilde B_j=W\hp B_j$, $\tilde C_j=C_jV$.  For us,
\[
\tilde D(s)=s\,(W\hp MV)-(W\hp AV)-k(s)\,(W\hp b)(c\hp V):
\]
\emph{the two large projected matrices are computed once; the design enters only through
the scalar $k(s)$ and the two small vectors $W\hp b$, $c\hp V$.}  This is the formal
reason for the design-independence measured below.

\emph{Interpolation (their Theorem 2.1, a special case of Beattie--Gugercin 2009).}  At a
point $\mu$ and for vectors $b_0,c_0$: if $D(\mu)^{-1}B(\mu)b_0\in\mathrm{Col}(V)$ then
$\tilde H(\mu)b_0=H(\mu)b_0$; if $(c_0\hp C(\mu)D(\mu)^{-1})\hp\in\mathrm{Col}(W)$ then
$c_0\hp\tilde H(\mu)=c_0\hp H(\mu)$; and if both, $c_0\hp\tilde H'(\mu)b_0=c_0\hp H'(\mu)b_0$.
Applied to the right/left singular vectors $v,w$ of $\sigma_1(H(\mu))$, this gives
\emph{Hermite interpolation of the largest singular value}: $\sigma_1$ and
$\partial_\omega\sigma_1$ of the reduced function equal those of the full function at
$\mu$.  The statement is \emph{per vector}, which matters for us below.  Two consequences
they draw: with only the input direction in $V$ one has $\sigma_1(H(\mu))\le\sigma_1(\tilde H(\mu))$
--- the two-sided reduced function is an \emph{interpolant}, not a bound; and $V$, $W$
must have equal dimension for $\tilde D$ to be square.

\emph{Algorithm 1 (greedy).}  Start with one interpolation point $i\omega_1$ and the
subspaces $V_1=D(i\omega_1)^{-1}B(i\omega_1)$, $W_1=(C(i\omega_1)D(i\omega_1)^{-1})\hp$.
Then repeat: compute $\omega_r=\arg\max_\omega\sigma_1(\tilde H_{r-1}(i\omega))$ by a
dense method on the reduced function; expand $V,W$ with the snapshots at $i\omega_r$;
orthonormalise.  Each iteration costs one sparse factorisation of $D(i\omega_r)$ and
$m+p$ solves.  \emph{Convergence (their Theorem 3.3)}: under a local convergence
assumption the iterates converge to a local maximiser of $\sigma_1$ \emph{superlinearly
in the subspace dimension}, and the proof uses only the Hermite interpolation at the last
two iterates; keeping all previous subspaces (their Remark 2.3) costs little and gives
much better global behaviour in practice, typically fewer than ten iterations to
$10^{-6}$.  \emph{Limitation (their \S4.4, stated by the authors)}: convergence is local
only, as for every large-scale $\mathcal L_\infty$ method to date; if the initial
interpolation points miss the part of the axis near a pole, the global maximiser is
missed --- their own example converges to $\omega=7.897$ when the maximiser is at
$1.092$.  Their remedy is more initial points; there is no certificate.  That limitation
is exactly what \S\ref{sec:cert} supplies, which is the point of \S\ref{sec:proj-mgmt}.

\subsection{Our variant: full-field channel, dominant directions, Woodbury columns}
\label{sec:proj-variant}
Our performance channel is not tall-and-skinny: $B=C\hp=M^{1/2}$ is $n\times n$, so
$D^{-1}B$ is the whole space.  Theorem 2.1 being per vector, the remedy is the
reduced-basis variant the authors themselves indicate: at each anchor $i\omega_a$ include
only the images of the \emph{dominant} singular directions,
\[
R\,M^{1/2}v_j\ \text{ in }V,\qquad R\hp M^{1/2}u_j\ \text{ in }W,\qquad j=1,\dots,n_{\mathrm{vec}},
\]
where $(u_j,\sigma_j,v_j)$ are the top singular triplets of $T(i\omega_a)$ --- which the
warm-started power/Lanczos iteration of \S\ref{sec:sigma} produces anyway --- plus the two
\emph{Woodbury columns} $R\,b$ and $R\hp c$, so that the rank-one controller correction
is representable inside the subspace exactly.  With $n_{\mathrm{vec}}\ge2$ the second
singular value is Hermite-interpolated as well.  Two projections are then in use:
the \emph{two-sided} one (interpolant, exact at anchors, no bound), for values,
$\sigma_2$, gradients and the eigenvalue predictor; and the \emph{one-sided input
restriction} $\|T\,V\|$ with $V$ orthonormal, which is a \emph{certified lower bound}
$\le\sigma_1$ at every frequency, since restricting the input can only lower the gain.
Gradients come from the master formula \eqref{eq:master} with the solved vectors replaced
by their subspace representatives $V\tilde v$, $W\tilde u$ --- no full solve.

\revQ{\emph{No $M^{1/2}$ is ever formed} (added October 5, 2026).  $M^{1/2}$ above is the
notation of the full-field channel ($B=C\hp=M^{1/2}$ on GL); on the backward-facing step the
channel is the weak-form one of \S\ref{sec:bfs}, $\sigma_1^2=\lambda_{\max}(Q\,T\hp Q\,T)$.
What the implementation stores per anchor is the closed-loop \emph{state response} $x_j$ and
\emph{adjoint response} $z_j$ of the leading singular pair (the vectors the Lanczos iteration
on $Q\,T\hp Q\,T$ produces, $x_j=T\,Q\,v_j$ up to scaling, $z_j$ its adjoint counterpart), plus
the Woodbury columns $R\,b$, $R\hp c$ and the images of the port bases.  The reduced singular
values are those of $G_V^{1/2}\tilde T\,G_W^{1/2}$ with $G_V=V\hp QV$, $G_W=W\hp QW$: the square
roots are of these $\tilde n\times\tilde n$ Gram matrices of the basis (a dense
$\tilde n\times\tilde n$ \code{sqrtm}), never of $Q$ or $M$, and no mass matrix is inverted
or factored.}

\subsection{What was measured on GL}
\label{sec:proj-numbers}
All tests dense, GL with $n=799$ (or $n=399$ where noted), performance channel
$T(\omega)=M^{1/2}R_{\mathrm{cl}}(\omega)M^{1/2}$, $n_{\mathrm{vec}}=3$, greedy anchors,
against a dense ground truth.

\emph{Plant sweep} ($\omega\in[-2,12]$, 141 points): anchors $1/2/4/5$ give maximum
relative error $7\cdot10^{-1}$ / $4.5\cdot10^{-1}$ / $7.8\cdot10^{-2}$ /
$1.7\cdot10^{-4}$ over the band; the peak is captured to $10^{-6}$ from the second anchor
on, and the greedy sequence chose the peak frequency second.  Five factorisations in
place of 141.

\emph{Closed loop at a fixed design} (Woodbury form, $r=8$ tridiagonal controller):
5 anchors, $\dim V=40$, maximum relative error $5.6\cdot10^{-6}$, peak to $2\cdot10^{-9}$.

\emph{Design variation, reusing $V$ built at $\theta_0$ without rebuilding} (controller
entries perturbed by relative amounts $\varepsilon$; actuator/sensor positions moved by
$0.15\varepsilon$ against Gaussian width $0.02$):
\begin{center}\small
\begin{tabular}{lcccc}\toprule
perturbation & peak drift & max rel.\ error $\sigma_1$ & peak error & lower bound holds\\
\midrule
controller, $\varepsilon=1.0$ & $68\%$ & $5.6\cdot10^{-6}$ & $6\cdot10^{-12}$ & yes\\
positions, $\varepsilon=0.5$ & $67\%$ & $5.6\cdot10^{-6}$ & $9\cdot10^{-7}$ & yes\\
positions, $\varepsilon=1.0$ & $37\%$ & $1.1\cdot10^{-4}$ & $1\cdot10^{-4}$ & yes\\
both, $\varepsilon=1.0$ & $53\%$ & $1.1\cdot10^{-4}$ & $1\cdot10^{-4}$ & yes\\
\bottomrule\end{tabular}
\end{center}
The accuracy is essentially independent of how far the design moves; the certified lower
bound never failed (worst shortfall $2\%$, a constant).  \emph{Plant-only subspace}
(built with no knowledge of $b$, $c$ or the controller, $\dim30$) reused on four very
different closed loops: $6\cdot10^{-3}$ in every case --- the subspace is a property of
the plant, and the two Woodbury columns buy three orders of magnitude.

\emph{$\sigma_2$} through the same subspace: $8.9\cdot10^{-6}$ (fixed design),
$3\cdot10^{-5}$ at $\varepsilon=1$.

\emph{Closed-loop eigenvalues} (pencil projected on $V$ augmented with the right/left
eigenvectors of the six rightmost eigenvalues, $\dim54$): the Ritz value nearest the true
rightmost eigenvalue tracks it to $10^{-13}$ at $\theta_0$, $5\cdot10^{-6}$ at
$\varepsilon=0.2$, $2\cdot10^{-5}$ at $\varepsilon=0.5$ --- an excellent
\emph{predictor} of where a watched eigenvalue went.  But the oblique projection also
produces \emph{spurious} Ritz values far in the right half-plane (one at $2.6+7.6i$,
corresponding to nothing): nothing variational protects a Petrov--Galerkin eigenvalue of a
non-normal pencil.  A full-model residual $\|(\Acl-\lambda\Mcl)v\|$ per Ritz pair, one
matvec, filters them.  \textbf{Predictor yes, certificate no.}

\emph{Real, unstable plant} (GL realified to the state $(\mathrm{Re}\,q,\mathrm{Im}\,q)$
--- real $A,M$ with GL's spectrum and its conjugate --- and shifted right until four
conjugate pairs are unstable, $N=798$): $\sigma_1$ on the axis is finite and symmetric;
a \emph{real} subspace built from the real and imaginary parts of the anchor columns,
with anchors on $\omega\ge0$ only (five: $0.4,0,1.6,4.0,0.8$; real $\dim60$), gives real
projected matrices and $9.6\cdot10^{-6}$ on \emph{both} halves of the axis, $\sigma_2$ to
$1.5\cdot10^{-5}$; the four unstable pairs' eigenspaces lie inside $V$ to a fraction
$1.000$ each without having been added, and the real reduced pencil reproduces them to
four digits, in conjugate pairs, with nothing spurious in the right half-plane.

\subsection{Unstable plants and real arithmetic}
The cavity plant is unstable and real; neither changes the recipe.  The anchors are solves
on the imaginary axis, where $s M-A$ is nonsingular because the plant eigenvalues are off
the axis; the resolvent near the axis is dominated by precisely the unstable modes, so the
greedy anchors capture them on their own (measured above) --- add their right and left
eigenvectors explicitly anyway, they are known from the global stability analysis, cost
nothing, and the left ones are needed for the eigenvalue predictor.  Stability of the
closed loop remains the certificate's statement (\S\ref{sec:winding}), not the reduced
model's.  Realness: the anchor factorisations stay complex (the shift is imaginary; a real
$2n$ reformulation gains nothing), but from the columns on everything is real --- split
each column into real and imaginary parts, the subspace is closed under conjugation, the
projected matrices are real, the reduced model is a real descriptor system with the
real-rational $k(s)$, its eigenvalues come in conjugate pairs, and anchors and sweep are
needed on $\omega\ge0$ only.  This composes with the realness savings of \S\ref{sec:ns}.

\revQ{\subsection{Optimise on the reduced model, enrich at the full one: Aliyev--Benner--Mengi--Voigt (2019), plus certification}
\label{sec:proj-mgmt}
The loop ``run the optimisation on the reduced model, then consult the full model at the
result and enrich there'' is not ours.  It is Algorithm~2.1 of N.~Aliyev, P.~Benner,
E.~Mengi, M.~Voigt, ``A subspace framework for $H_\infty$-norm minimization'',
\emph{SIAM J.\ Matrix Anal.\ Appl.}\ 41(2) (2020) 928--956 (arXiv:1905.04208, 2019),
written for exactly our problem: minimise $\|H[\mu]\|_{H_\infty}$ over the parameters $\mu$
of a large descriptor system whose matrices depend affinely on $\mu$ through scalar
functions, $E(\mu)=\sum f_j(\mu)E_j$, $A(\mu)=\sum g_j(\mu)A_j$, $B(\mu)=\sum h_j(\mu)B_j$,
$C(\mu)=\sum k_j(\mu)C_j$ --- which ours do: the controller entries through the rank-one
coupling terms, the field coefficients through $b(\beta)=\Sigma\beta$.  The algorithm:
\begin{enumerate}[leftmargin=1.8em,itemsep=1pt]
\item $\mu^{(k)}\leftarrow\arg\min_\mu\|H^{V,W}[\mu]\|_{H_\infty}$ --- the reduced norm
minimised over the design (they use GRANSO, or the global optimiser \code{eigopt} when
$\mu$ has few components);
\item $\omega^{(k)}\leftarrow\arg\max_\omega\sigma_1(H[\mu^{(k)}](i\omega))$ --- one
\emph{full-model} norm evaluation at the reduced optimum;
\item $V\leftarrow V\oplus D(\mu^{(k)},i\omega^{(k)})^{-1}B(\mu^{(k)})$,
$W\leftarrow W\oplus D(\mu^{(k)},i\omega^{(k)})^{-\ast}C(\mu^{(k)})^\ast$; repeat.
\end{enumerate}
Their Lemma~2.1: at every iterate the reduced function Hermite-interpolates the full one
\emph{in both $\omega$ and $\mu$} --- value, $\sigma_2$, $\partial_\omega\sigma=0$, and
$\nabla_\mu\|H\|_{H_\infty}=\nabla_\mu\sigma^{V,W}$ (the reduced gradient with respect to
the design equals the true gradient at the design where the snapshot was taken).  Their
Theorem~3.8: \emph{superlinear} convergence of $\mu^{(k)}$ to a minimiser, under
Assumption~3.1 --- the supremum of $\sigma_1(\mu_\ast,\cdot)$ is attained at a \emph{unique}
$\omega_\ast$ where $\sigma_1$ is simple --- and with $\mu_\ast$ interior (box constraints
active at $\mu_\ast$ are handled by restricting to the inactive variables, their
Remark~3.9; their real examples all have optima on the box boundary).  In practice
$\le10$ outer iterations and ``a few seconds for order $10^4$''.  Stability is explicitly
left open: ``efficient stability checks would be needed for controller design''.

\emph{What is ours.}  Two things, and both are consequences of the kind of optimum we
meet.  (i)~Their Assumption~3.1 --- a single active peak --- fails at our optima, which
are balances of several peaks (\S1 of the optimisation note); at such a point the
reduced function must interpolate the full one at \emph{every} active peak, or the
reduced optimum drifts along the crease.  We therefore enrich at all certified
near-active frequencies --- the working set $\Omega$ --- rather than at the single
maximiser of line~2; their Algorithm~2.2, which adds nearby interpolation points to
control second derivatives, points the same way.  (ii)~Line~2 trusts a full-model norm
evaluation; ours is the certified sweep of \S\ref{sec:cert}, and the certificate does
double duty: it decides whether the reduced level is the true level, and \emph{where} to
enrich --- a refuted point, or a certified peak whose reduced value disagrees beyond
tolerance, is the next anchor.  Superlinear locally (theirs), globally certified (ours),
and the reduced model is never trusted.  Their per-design snapshot (line~3 at the
current $\mu$) is our per-design refresh of the Woodbury columns; for field parameters
the affine dependence $b(\beta)=\Sigma\beta$ can be projected once, $W\hp\Sigma$ and
$\Sigma\hp V$ ($\dim V\times n_\beta$), so that the reduced model carries the
$\beta$-dependence of $B$ exactly and the refresh only sharpens the resolvent images.
Two practical gifts of their setting: on a reduced model an evaluation costs microseconds,
so random multistart of the inner problem is free, and for designs with few parameters
(the position cases, $n_\theta=2$--$4$) the reduced problem can be optimised
\emph{globally} (\code{eigopt}: piecewise-quadratic support functions, Breiman \& Cutler
1993) before certification.  What the certification cannot make cheaper is itself; the
reduced model shortens its bisection and predicts the hot windows.}

\revQ{\subsection{Eigenvalues and the margin on the reduced model: Mengi (2022, 2023)}
\label{sec:proj-mengi}
Two further papers by the same group supply the mechanisms this section had left as
``predictor only'' and ``certificate outside''.

\emph{Watch-list tracking by interpolation.}  E.~Mengi, ``Large-scale estimation of
dominant poles of a transfer function by an interpolatory framework'', \emph{SIAM J.\
Sci.\ Comput.}\ 44(4) (2022) A2412--A2438: compute the dominant poles of the
\emph{reduced} system (QZ on the small pencil), then expand $V,W$ so that the full
system is Hermite-interpolated \emph{at those poles} --- solves with $(A-\lambda E)$ at the
complex shift $\lambda$ equal to the pole estimate, with $q\ge1$ derivative directions
$(A-\lambda E)^{-1}E(A-\lambda E)^{-1}B$ --- and repeat; convergence at least quadratic
(his Theorem~4.3), termination and per-pole convergence by the full-model residual
$\|(A-\lambda E)Vv\|_\infty$ (his \S5.2--5.3), one LU per pole per iteration, more
reliable than SAMDP on benchmarks.  For us this replaces the continuation gate's
segment-by-segment tracking: the watch-list eigenvalues are followed through the reduced
pencil as $\theta$ changes (the predictor of \S\ref{sec:proj-numbers}, $5\cdot10^{-6}$),
each estimate is polished by one shift-invert solve at the Ritz value --- his
interpolation step --- and the residual is the convergence certificate for that pole.
The eigenvalue-identity problem that broke endpoint checks disappears, because the
reduced pencil's poles vary continuously with $\theta$; the spurious Ritz values of
\S\ref{sec:proj-numbers} are handled by the same step --- interpolating at a non-pole adds
information, and the value either converges to a true pole or fails the residual test.
His dominance measure $\|Cv\|\,\|w\hp B\|/|\mathrm{Re}\,\lambda|$ becomes, for our
full-field channel $B=C\hp=M^{1/2}$, $1/|\mathrm{Re}\,\lambda|$ times the mode's
biorthogonality factor: the dominant poles are the near-axis, non-normal modes --- the
watch list and the peak generators at once, one subspace for both.  His initialisation
(\S5.1) is ``run the $L_\infty$ subspace method crudely'', i.e.\ our anchors.  The winding
count of \S\ref{sec:winding} stays as the unconditional statement; this mechanism makes
the per-step gate rigorous per pole and cheap.

\emph{The margin as a constraint on the reduced model.}  E.~Mengi, ``A subspace framework
for $L_\infty$ model reduction'', arXiv:2309.08011 (2023), \S6: the reduced problem is
minimised subject to $\alpha(A_{\mathrm{red}},E_{\mathrm{red}})\le\beta$, the spectral
abscissa of the \emph{reduced} pencil, imposed by a log-barrier
$L_\mu=F-\mu\log(\beta-\alpha)$ --- and he shows that the Hermite interpolation
properties between full and reduced objective \emph{extend to the barrier function},
gradient and Hessian included.  For us that is the cheap form of the structural fix the
optimisation note specifies (its to-do item~2): the explicit margin constraint
$\alpha(\theta)\le-\eps$, evaluated and differentiated on the reduced closed-loop pencil
inside the inner loop --- eigenvalues by QZ on a $\dim V$ pencil, gradients from the
reduced eigenvectors, residual filter applied --- and certified afterwards by the
winding count.  The same paper's diagnosis of direct smooth optimisation of an $L_\infty$
objective --- ``at best a linear rate, requiring too many evaluations of the costly
norm'' --- is the v2 direct-mode record of the optimisation note in one sentence.}

\subsection{What projection does not replace: the frozen-preconditioner test}
\label{sec:proj-precond}
A related idea --- reuse one factorisation as a \emph{preconditioner} for GMRES at
neighbouring frequencies --- was tested first and rejected for the $\sigma_1$ oracle, for
two measured reasons.  The preconditioned operator $I+i\Delta\omega\,R(\omega_0)M$ has
eigenvalues $1+i\Delta\omega/(i\omega_0-\lambda_j)$: on GL's dense branch there are $112$
plant eigenvalues within $2r_0$ of a shift at $\omega=8$, so the spectrum spreads instead
of clustering (266 of 799 eigenvalues further than $0.1$ from $1$ at $\Delta\omega=r_0$),
and non-normality widens the gap between $\|R(\omega_0)M\|$ and $1/r_0$ by a factor $18$
at $\omega=0$.  GMRES then needs $10$--$30$ iterations at tolerance $10^{-6}$; and since
the oracle needs $m\approx10$--$20$ right-hand sides per frequency, each paying its own
iterations, break-even is $k^\ast=1+C/m$ with $C$ the LU-to-solve cost ratio ---
$k^\ast\approx4$ in 2-D --- and the idea loses by $3$--$7\times$.  It \emph{wins} where
$m=1$: the winding certificate's samples ($G=c\hp Rb$, tolerance $10^{-3}$), measured
at $3$--$16$ iterations for $\Delta\omega$ up to two band widths, against $k^\ast\approx50$.
Projection replaces the oracle; anchor-preconditioning serves the winding tail; the two
are complementary, and both hang on the same handful of resident factorisations.

\subsection{Cost and parallelism at cavity scale}
\label{sec:proj-cost}
For Taylor--Hood P2--P2--P1 on $2\cdot10^5$ triangles ($7$--$9\cdot10^5$ dof), a complex
sparse LU with nested dissection should run to $50$--$100$ nonzeros per row in the factors:
\emph{estimates}, to be measured, of $1$--$3$\,GB and $30$--$120$\,s per factorisation on
one node.  Fifty of those per inner iteration, times hundreds of iterations, times the
case count, is the cluster; with projection the accounting changes shape: (i)~anchors,
$5$--$10$ factorisations \emph{once per mesh}, design-independent, factors stored
($\sim20$\,GB) --- parallel \emph{across} anchors, multithreaded within one node, never
distributed across nodes for a 2-D factorisation; (ii)~inner loop, \emph{zero}
factorisations --- per design two solves per anchor (field cases only), per frequency
$\sim\dim V$ sparse matvecs and a $\dim V\times\dim V$ SVD, all of $\Omega$ on one node;
(iii)~certification, $O(10^2)$ factorisations per outer round, embarrassingly parallel
over shifts.  The node estimate goes from ``fifty nodes held for every case'' to ``one
fat node for the optimisation and a handful borrowed for minutes per outer round''.
The open number is the subspace dimension in 2-D: GL needs $30$--$60$; the cavity's richer
near-axis spectrum may need several times more, which is still small but decides the
estimate above, and it is the first thing to measure.}

%======================================================================
\section{The backward-facing step: algebra without $M^{-1}$}\label{sec:bfs}

{\small [Added October 2, 2026 (branch \code{bfs}; implementation notes in \code{BFS.md}).
The plant is the rounded backward-facing step of Dergham, Sipp \& Robinet (2013) at
$Re=600$, Taylor--Hood P2--P2--P1, $n=191{,}720$; the devices are those of Barbagallo,
Dergham, Sipp, Schmid \& Robinet (JFM 2012) with free positions.  Rule of this section:
\emph{no inverse and no Cholesky factor of the velocity mass}; $M$ is only multiplied,
scaled and added.]}

\subsection{Weights in weak form: $BB\tp=C\tp C=Q$}
The state $q=(u,p)\in\R^n$; $Q=M$ is the mass matrix, zero on the pressure block and on the
Dirichlet rows and columns; $Q'$ is its velocity block ($n_v\times n_v$, positive definite on
the free velocity dofs); $P$ is the prolongation $\R^{n_v}\to\R^n$ that pads a zero pressure,
so $Q=PQ'P\tp$.  The disturbance $w$ is a velocity field with $\|w\|^2=w\tp Q'w$ entering in
weak form, $Q\dot q=Aq+QPw+b\,u$; the performance is the whole state in the $Q$ seminorm.  With
the normalised input $w=Q'^{-1/2}\hat w$ and output $z=Rq$, $R\tp R=Q$:
\[
B=QPQ'^{-1/2}\ (n\times n_v),\qquad
BB\tp=QPQ'^{-1}P\tp Q=PQ'Q'^{-1}Q'P\tp=Q,\qquad
C\tp C=R\tp R=Q .
\]
$B$ and $C\tp$ have different shapes, but their Gram products are the same $n\times n$
matrix, and the Hamiltonian/even pencil needs only these products.  $Q'^{-1/2}$ and $R$
appear in the derivation and cancel; they are never formed.  (For GL the input entered in
strong form, $E\dot q=Aq+E^{-1/2}w$, giving $BB\tp=E^{-1}$; premultiplying by
$G=\mathrm{diag}(E,I,\dots)$ removes it.  In weak form the inverse never appears.)  Hence
\[
\sigma_1^2(\omega)=\lambda_{\max}\bigl(Q\,T\hp Q\,T\bigr),\qquad T=(i\omega M_{\mathrm{cl}}-A_{\mathrm{cl}})^{-1},
\]
an operator self-adjoint in the $Q$ semi-inner product (real nonnegative spectrum), and
the $\gamma$-pencil of \S\ref{sec:cert} is
\[
\begin{pmatrix}A_{\mathrm{cl}}&\gamma^{-2}Q_{\mathrm{cl}}\\ Q_{\mathrm{cl}}&A_{\mathrm{cl}}\hp\end{pmatrix}
-s\begin{pmatrix}M_{\mathrm{cl}}&0\\0&-M_{\mathrm{cl}}\end{pmatrix},
\qquad Q_{\mathrm{cl}}=\mathrm{diag}(Q,0_r),
\]
sparse plus the rank-4 controller coupling (Woodbury), exactly as on GL.  The archive's
optimal gains are reproduced to $10^{-14}$--$10^{-15}$ with it.

\subsection{Devices}
Actuator: a vertical body force $\psi_a=(0,g)$,
$g=\exp(-|x-x_a|^2/2\sigma^2)$, $\sigma=0.5/\sqrt{\ln2}$, load $b=M\psi_a$, free centre
$x_a\in\R^2$.  Sensor: the wall shear on the floor $y=0$,
$y=\int\chi(x-x_s)\,\partial_yu\,dx$, i.e.\ the vector $c=S\tp(w\odot\chi)$ where
$(Sq)_k=\partial_yu(x_k,0)$ at Gauss points $x_k$ (built once from the P2 basis) and $\chi$ a
length-$0.6$ window with $\tanh$ edges of width $0.05$; free start $x_s\in\R$.  Both are
unnormalised.  $\partial b/\partial x_a=M\,\partial\psi_a/\partial x_a$ and
$\partial c/\partial x_s=S\tp(w\odot\partial_{x_s}\chi)$, so the placement gradients are the
GL formulas with these vectors (checked against finite differences to $10^{-9}$).  The
sensor is \emph{not} of the form $M\psi$: this matters for the tail bound below.

\subsection{Eigenvalue-free disks without $M^{-1}$, and why they do not help here}
For a pencil $(\mathcal A,\mathcal E)$ and a shift $s$, an eigenpair with $\mathcal Ev\ne0$
gives $v=(\lambda-s)Xv$, $X=(\mathcal A-s\mathcal E)^{-1}\mathcal E$, so
$|\lambda-s|\ge1/\rho(X)$.  Write $\mathcal E=R\tp JR$ with $J$ a signature and
$W=R\tp R\succeq0$ ($\mathcal E_H=\mathrm{diag}(M_{\mathrm{cl}},-M_{\mathrm{cl}})$: $W=\mathrm{diag}(M_{\mathrm{cl}},M_{\mathrm{cl}})$).  The nonzero
spectrum of $X=P^{-1}R\tp JR$ is that of $RP^{-1}R\tp J$, so
\[
\rho(X)\le\|RP^{-1}R\tp\|_2,\qquad \|RP^{-1}R\tp\|_2^2=\lambda_{\max}\bigl(P^{-\mathsf H}WP^{-1}W\bigr),
\]
the second equality again by cycling.  The operator on the right is self-adjoint in the
$W$ semi-inner product: Arnoldi on the factors of the shift, no inverse of $W$, $R$ never
formed.  For an SPD mass this equals the earlier $\lambda_{\max}(X\hp WX,W)$ (similar
operators, $JW=WJ$); the code now uses one form for both plants (equal to $10^{-15}$ on GL).

\emph{Measured on the step} (test plant, a stabilising design, $\gamma=1645$): the free radius
at $i\theta$ is $10^{-7}$--$10^{-10}$ where the nearest eigenvalue of the even pencil is
$0.015$--$0.5$ away.  The convective resolvent is so non-normal that norm-based radii are
useless; the disks of the $\gamma$-certificate must come from the \emph{nearest eigenvalues}
(shift-invert Arnoldi, converged), the Schwerdtner--Mengi--Voigt construction.  The free
radius is therefore switched off for the step, and an Arnoldi set that does not converge is
retried with a larger subspace or the interval bisected.

\subsection{The numerical range on divergence-free velocities, certified by inertia}
$A=\begin{pmatrix}A_{vv}&B\tp\\B&0\end{pmatrix}$ (Dirichlet rows $-I$), $M=\mathrm{diag}(M_{vv},0)$.
An eigenvector $(u,p)$ has $Bu=0$, $u=0$ on Dirichlet dofs, and
$\lambda=u\hp A_{vv}u/u\hp M_{vv}u$ ($u\hp B\tp p=(Bu)\hp p=0$): the spectrum lies in the
numerical range $W$ of $A_{vv}$ in the $M$ inner product on $\ker B$.  Its bounds are tops of
Hermitian pencils $(H,M)$ on $\ker B$: $H=\tfrac12(A+A\tp)$ for $\max\mathrm{Re}$, $\pm\tfrac1{2i}(A-A\tp)$
for the heights, and the S-lemma combinations $\pm S+\tau(H-s_0M)$ of \S\ref{sec:ns}.
Certificate: by the saddle-point inertia theorem,
\[
\sigma M-H\succ0\ \text{on}\ \ker B
\iff
\mathrm{In}\begin{pmatrix}\sigma M_{vv}-H_{vv}&B\tp\\B&-\delta I\end{pmatrix}=(n_v,\,n_p,\,0)
\]
for small $\delta>0$, and the implication ``$\Leftarrow$'' holds for \emph{every} $\delta>0$ (on
$\ker B$ the term $B\tp B/\delta$ vanishes) --- the direction a certificate needs.  The
inertia is read from one sparse $LDL\tp$ without pivoting (the matrix is quasi-definite
exactly when the test passes; a breakdown or a wrong count answers ``not certified''), by
Sylvester's law.  A complex Hermitian $H_r+iH_i$ is embedded as the real symmetric
$\begin{pmatrix}H_r&-H_i\\H_i&H_r\end{pmatrix}$.  $\lambda_{\max}\le\sigma$ is then found by doubling and
bisection on $\sigma$: every accepted $\sigma$ is a proof, no Ritz value from below.  This is
an $LDL\tp$ of $\sigma M-H$, never of $M$.  (MUMPS reports no inertia in complex arithmetic,
so the factorisation is \code{qdldl}'s, in real arithmetic: $3$\,s at $n=191{,}720$.)
Validated on a $1{,}279$-dof mesh against dense generalised eigenproblems on $\ker B$:
equal to $10^{-7}$, from above.  On the test plant $\max\mathrm{Re}\,W\le1.76$, $|\mathrm{Im}|\le87.9$.

\subsection{The winding count for a wall-shear sensor: an $H^1$ tail and an $L^2$ tail}
The count of \eqref{eq:count} needs $|k(s)G(s)|<\tfrac12$ beyond a frequency $W_T$, from a
bound on $G(s)=c\hp(sM-A)^{-1}M\psi_a$.  With $x=(u,p)$ the response,
\[
s\|u\|_M^2-u\hp A_{vv}u=(u,\psi_a)_M,\qquad Bu=0 .
\]
(i) $\mathrm{dist}(s,W)\,\|u\|_M^2\le\|u\|_M\|\psi_a\|_M$, so $\|u\|_M\le\|\psi_a\|_M/d$,
$d\le\mathrm{dist}(s,W)$ from the heights (\code{numrange\_distance}).  (ii) The shear sensor
is not bounded in $L^2$ in the continuum (a trace of the gradient; the discrete one is, with a
constant that grows as the wall mesh is refined: (iv) below), but it is in $H^1$:
$|c\hp u|\le\|c\|_{L^{-1}}\|u\|_L$, $L$ the velocity stiffness,
$\|c\|_{L^{-1}}^2=c\tp L^{-1}c$ (one solve).  (iii) With
$\beta\ge\lambda_{\max}(\tfrac12(A_{vv}+A_{vv}\tp)+\nu'L,\,M)$ on $\ker B$ (certified as above,
$\nu'=\nu/2$), $-\mathrm{Re}\,u\hp A_{vv}u\ge\nu'\|u\|_L^2-\beta\|u\|_M^2$, and the real part of
the identity gives $\nu'\|u\|_L^2\le\|u\|_M\|\psi_a\|_M+(\beta-\mathrm{Re}\,s)_+\|u\|_M^2$.  Hence
\[
|G(s)|\le\|c\|_{L^{-1}}\|\psi_a\|_M\sqrt{\frac{1/d+(\beta-\mathrm{Re}\,s)_+/d^2}{\nu'}}
=O(d^{-1/2}).
\]
(iv) An $L^2$ bound from (i) alone (2026-10-16): $c$ is zero off the free velocity dofs,
so Cauchy--Schwarz in $M_{vv}$ (the velocity mass on them, SPD) gives
$|c\hp u|=|(M_{vv}^{-1}c_v)\hp M_{vv}u|\le\|c\|_{M_{vv}^{-1}}\|u\|_M$ and
\[
|G(s)|\le\|c\|_{M_{vv}^{-1}}\|\psi_a\|_M/d=O(d^{-1}),
\]
with $\|C\|_{M_{vv}^{-1}}^2=\lambda_{\max}(C_v\hp M_{vv}^{-1}C_v)$ for several sensors.  It is
computed without factorising or inverting $M_{vv}$: CG gives $Y\approx M_{vv}^{-1}C_v$, and with
$R=C_v-M_{vv}Y$ the identity $C_v\hp M_{vv}^{-1}C_v=C_v\hp Y+Y\hp R+R\hp M_{vv}^{-1}R$ bounds it by
$\lambda_{\max}(C_v\hp Y+Y\hp R)+\|R\|_2^2/\lambda_{\mathrm{floor}}$, $\lambda_{\mathrm{floor}}$ the
element floor $\min_e\lambda_{\min}(M_e)\le\lambda_{\min}(M_{vv})$ (thrift level (b) of the GL tail,
\S\ref{sec:winding-num}; interlacing: $M_{vv}$ is a principal submatrix of the assembled mass,
and the floor's premise is checked by reassembling $M$ from the mesh).  The count uses
$\min(H^1, L^2)$.  Both bounds read $c\hp x$ through the velocity alone: $c$ must be zero on the
pressure dofs (else no tail bound at all, the count's \code{TailBoundError}), while entries on
the Dirichlet dofs, where the response is zero, leave the $H^1$ bound (the $L^2$ one keeps its
stricter premise).  Test plant, decentral $2\times2$ record $369.87$: $\|C\|_{M_{vv}^{-1}}=738.05$
(floor $1.94\times10^{-5}$ against $\lambda_{\min}=1.08\times10^{-4}$); the $L^2$ bound is the
smaller from $w\approx100$ ($51.9$ against $102.5$), and $W_T$ falls from $7.86\times10^6$ to
$7680$.  ERRATA (2026-10-16): this subsection claimed the $d^{-1/2}$ decay was genuine, from the
Stokes layer of thickness $\sqrt{\nu/\omega}$ and its wall shear $(\nu\omega)^{-1/2}$.  That is a
continuum argument.  The discrete plant decays like $1/w$ ($w\|G(-0.05+iw)\|=11.74$ from $10^2$
to $10^7$ at that record): the mesh does not resolve a layer that thin, and the crossover moves
out on a finer wall mesh.  With the $H^1$ bound alone, $W_T\sim(2|K|\,C)^2$ grew with the square of
the feedthrough and reached the count's old cap of $10^7$ in 20261010n.  With $(iv)$ it grows
linearly in it, as for an $M\psi$ sensor.  $P_{\mathrm{plant}}$ and the near-line
plant poles come from the enclosure of \S\ref{sec:ns} with the inertia bounds above and the
tiles' nearest-eigenvalue exhaustion; they are design-independent and cached per mesh.
The plant is real, so $f(\bar s)=\overline{f(s)}$ and the march covers $[0,W_T]$ only, the phase
change doubled.  Measured on the test plant ($n=33{,}593$; $\max\mathrm{Re}\,W\le1.758$,
$\beta=2.48$; no plant pole right of $-0.025$, $282$ within $0.6$ of it) at a design with
$\alpha=-0.0501$, $K=0.0046$: $W_T=120$, $3617$ points, winding $-4.6\times10^{-15}$, count $0$.

\subsection{The abscissa, and random starts}
The line certificate of \S\ref{sec:cert} proves only that nothing lies \emph{on}
$\mathrm{Re}\,s=-\eps$; it reports the eigenvalues near its shifts.  For an unstable design it
is not the abscissa.  Measured: a HIFOO phase driven by it at $\mathrm{Re}=-0.08$ stopped at
``$\alpha=-0.0655$'' while the certificate at $-0.05$ then found $-0.0438+0.66i$.  The count
above gives the abscissa properly: $\#\{\mathrm{Re}\,\lambda>\ell\}$ at any $\ell$, bisected on
$\ell$.  For the starts in use, feasibility is reached by continuation in the loop gain
$\tau(n,K)$ from the open loop (stable plant, compensator shifted into the margin), every
step gated like an accepted design; the count certifies the result independently of the
chain.

\section{Cost accounting}\label{sec:cost}

\begin{center}
\small
\begin{tabular}{lll}
\toprule
quantity & primitives per evaluation & notes\\
\midrule
$\sigma_1(\omega)$ & 1 LU $+\ O(10)$ solves & warm-started power iteration \\
$\nabla_\theta\sigma_1(\omega)$ (\emph{all} params) & $+\,2$ solve sets & independent of NZ \\
$\partial_\omega\sigma_1$, Newton peak polish & inner products $+\ O(3)$ values & per peak \\
certified band maximum & $O(10^2)$ shift-invert solves & rtol $10^{-6}$, adaptive disks \\
$\alpha$ tracking (per accepted step) & few bordered solves per eigenvalue & shifted at previous location \\
winding certificate (per saved design) & $500$--$900$ LUs of $R(s)^{-1}$ & unconditional count \eqref{eq:count} \\
\revP{reduced sweep on $\Omega$ [proposed, \S\ref{sec:proj}]} & \revP{$0$ LU; $\sim\dim V$ matvecs $+$ one $\dim V$ SVD per frequency} & \revP{anchors: $5$--$10$ LUs once per mesh} \\
$\nabla_\theta\alpha$ & one eigenpair ($+$ adjoint) & formula \eqref{eq:alphagrad} \\
\bottomrule
\end{tabular}
\end{center}

\noindent
Every line is (P1)--(P3).  The two design decisions that make the whole pipeline scale are
visible in this table: gradients whose cost does not grow with the number of parameters
(\S\ref{sec:grads}), and an outer method (the exchange) shaped so that the LU-bearing
lines are executed for many independent frequencies concurrently.  The controller only
ever appears through the scalar $k(s)$ of \eqref{eq:ks} and the rank-2 identity
\eqref{eq:sherman} --- which is also, through \eqref{eq:factorised}, the reason the
stability of an $(n{+}r)$-dimensional loop can be certified by winding a single complex
number around a line.

\paragraph{Parallel cost, corrected.}  The winding line of the table is a \emph{serial}
count.  Split at the knots (\S\ref{sec:winding-num}), the same work
is one round of $\sim120$ knot factorisations followed by a work queue whose critical
path is the longest interval --- a few tens of serial factorisations --- so the wall time
at cavity scale is set by the hottest pole window, not by the total.

\section{Decisions and pending items}
\label{sec:todo}

The review rounds produced decisions and open items; they are collected here, each
pointing at the section that justifies it.  Items 1--2 amend statements of earlier
editions; the rest are code.
\begin{enumerate}[leftmargin=1.8em,itemsep=2pt]
\item \textbf{Tail bound for a singular mass matrix} (\S\ref{sec:winding-num}).  The bound $\|R(s)\|\le1/(\lambda_{\min}(M)(|s|-\nu))$ presumes
$M$ invertible; the cavity's $M=\mathrm{diag}(M_v,0)$ is not.  Re-derive on the
discretely divergence-free velocity subspace with the Leray projector $\Pi$:
$\nu:=\|\Pi M_v^{-1}A_{vv}\Pi\|$, $\lambda_{\min}(M_v)$, and a strictly proper $G$
guaranteed by momentum-equation inputs and velocity outputs.  Record the two modelling
rules that fall out (pressure sensors make $G$ proper, not strictly proper, and shift the
tail disk to $1-KG(\infty)$; inputs must never enter the constraint equation).  Validate
on a Stokes or lid-driven benchmark where a dense generalised eigensolve still exists.
\emph{Done for the step} (\S\ref{sec:bfs}): the range on $\ker B$ by inertia, with no
projector and no $M_v^{-1}$, and an $H^1$ tail for the wall-shear sensor; validated on a
$1{,}279$-dof mesh against dense eigensolves on $\ker B$.
\item \textbf{Once-per-mesh constants, tight values with certified bounds}.  Replace the ARPACK \code{which='SM'} call in
\code{plant\_tail\_constants} by shift-invert Lanczos (Cholesky- or CG-based) for a
\emph{tight} $\lambda_{\min}(M)$, and compute $\|A\|_2$ by Lanczos on $A\hp A$ (matvecs
only); alongside, compute the element-wise floor
$\lambda_{\min}(M)\ge\min_e\lambda_{\min}(M_e)$ and the H\"older ceiling as certified
cross-checks.  Tight constants shorten the tracked segment on every design; the bounds
catch bugs in the tight computations and vice versa.
\item \textbf{$P_{\mathrm{plant}}$ at cavity scale} (\S\ref{sec:contour}).
Implement the three steps as code, not prose: numerical-range rectangle with the inertia
certificate (one Hermitian Lanczos plus one $LDL\hp$ per bound); Beyn contour count with
probes, cross-validated at two quadrature orders and two probe blocks, left edge at
$\ell-\delta$; shift-invert polish of every counted eigenvalue.  Engine decision:
\textbf{Beyn}, per \S\ref{sec:contour}.  Prototype at GL against the dense count that
\code{plant\_pole\_data} already provides.  \emph{For the step} (\S\ref{sec:bfs}) the rectangle
uses the inertia certificate as here, with a sparse $LDL\tp$ in place of Lanczos (bisection on
$\sigma$), and the tiles are exhausted by nearest-$k$ shift-invert instead of Beyn; Beyn
remains the engine to decide for the cavity.
\item \textbf{Parallel winding certificate} (\S\ref{sec:winding-num}).
Interval decomposition at the knots, two-pass (knot values, then a work queue with
midpoint splitting of intervals that bisect below a threshold), phase subtotals summed.
\code{winding\_cert.py} is serial as frozen; measure the critical path at GL before
projecting cavity wall times.
\item \textbf{A claim in the code that the code does not honour --- resolved: implement,
merged into item 4.}  The docstring of \code{winding\_cert.py} advertises a
\code{--rigor} mode (a Neumann segment bound per sample, ``$\sim8$ extra matvecs per
point''); the argument parser has no such option and no such bound is implemented.
What was meant is a mechanism of this note --- the per-interval Neumann bound on $|kG|$
that closes quiet intervals without further samples, i.e.\ exactly the
``coarse-certify'' pass-2 test of the parallel scheme (item~4) and thrift item~(4) of
the cavity-cost answer.  So: implement it as part of item~4's pass~2, where it belongs;
the docstring then becomes true.  Recorded in the spirit of the report's \S7.
\item \textbf{Cost table} (\S10).  Add a column for the parallel critical path once
item~4 is measured, so the table stops quoting serial counts for a quantity that will
never run serially at scale.
\item \textbf{MIMO capability} (\S\ref{sec:structure}).  Generalise the
coupling to $n_a$ actuators and $n_s$ sensors: Woodbury in place of Sherman--Morrison,
the return-difference determinant $f=\det(I-\mathsf K\mathsf G)$ in place of the scalar,
$n_a{+}n_s$ solves per factorisation in place of two.  Structure already in place; cost
grows only in the solve count.
\item \textbf{Projection oracle, GL prototype} (\S\ref{sec:proj}).  Implement the
reduced-basis variant in the sparse pipeline: greedy anchors on $\omega\ge0$, real
subspace, top-3 singular directions plus the two Woodbury columns per anchor, unstable
eigenvectors added; values, $\sigma_2$, gradients and the eigenvalue predictor from the
reduced model; certification on the full model with refuted points and certified-peak
disagreements as the enrichment rule \revQ{--- Algorithm~2.1 of Aliyev--Benner--Mengi--Voigt
with enrichment at $\Omega$ instead of the single maximiser and the certified sweep in
place of line~2 (\S\ref{sec:proj-mgmt}); project the field basis once ($W\hp\Sigma$,
$\Sigma\hp V$) so that $b(\beta)$ is exact in the reduced model.}  Validate against the 27 frozen cases: same certified
optima, factorisation count per inner iteration reduced from $|\Omega|$ to zero.  Then
the \emph{trajectory} test that has not been run: does $V$ ever need refreshing along a
full optimisation, and does the certified-vs-reduced gap at the peaks detect it.
\revP{\item \textbf{Subspace dimension in 2-D --- measure first.}  On the cavity mesh
(or a coarser cavity), build the plant subspace once and record the greedy dimension
needed for $10^{-5}$ over the band; every cost estimate of \S\ref{sec:proj-cost} depends
on it.  Store the anchor factorisations and measure their size and time
(estimates: $1$--$3$\,GB, $30$--$120$\,s).}
\revP{\item \textbf{Anchor-preconditioned winding tail} (\S\ref{sec:proj-precond}).
GMRES from the nearest resident anchor for the single-solve samples of the segment,
tolerance $10^{-3}$, with the frozen $|f|$ used to tighten the tolerance near closed-loop
poles; the $\sigma_1$ oracle and the certified disk radii explicitly excluded until the
inexact-shift-invert bound is written.  Merge with item~4's two-pass scheme.}
\item \textbf{Watch-list tracking by interpolation} (\S\ref{sec:proj-mengi}).  Replace
the continuation gate's tracking by Mengi's dominant-pole framework: reduced-pencil poles
as $\theta$ moves, one shift-invert polish at each Ritz value, full-model residual as the
per-pole convergence test; spurious Ritz values eliminated by the same step.  Validate by
replaying the 27 frozen accepted-step chains against the gate's decisions.
\revQ{\item \textbf{Margin constraint on the reduced model} (\S\ref{sec:proj-mengi}).
Impose $\alpha_{\mathrm{red}}(\theta)\le-\eps$ inside the inner loop --- QZ on the reduced
closed-loop pencil, gradients from reduced eigenvectors, residual filter --- as the cheap
form of the optimisation note's item~2; winding count outside.  Test on K+B+C, the case
that stalled at the boundary.}
\revQ{\item \textbf{Global optimisation of the reduced problem for the position cases}
(\S\ref{sec:proj-mgmt}).  With $n_\theta=2$--$4$, run a global optimiser
(\code{eigopt}-style) on the reduced model, then certify; and use free multistart of the
reduced inner problem in the random-start study of the optimisation note's item~1.}
\end{enumerate}

\end{document}
